\documentclass[11pt]{article}

\usepackage[margin=0.8in]{geometry}
\usepackage[T1]{fontenc}
\usepackage{lmodern}
\usepackage{microtype}
\usepackage{graphicx}
\usepackage{booktabs}
\usepackage{array}
\usepackage{amsmath,amssymb,tabularx}
\usepackage{xcolor}
\usepackage[font=small,labelfont=bf]{caption}
\usepackage{placeins} % floats stay in their section
\usepackage{needspace} % keeps short headings with what follows
\usepackage{xurl}  % long commit URLs break across lines
\usepackage[hidelinks]{hyperref}
% Compact contents: sections, references and appendices (no subsections), one line each, dotted leaders to the
% page number at the right margin, little space between entries. The appendix writes \appendixcontents into the
% .toc, which prints an "Appendices" label before the appendix entries.
\providecommand{\appendixcontents}{}
\makeatletter
\newcommand{\compactcontents}{\begingroup\setlength{\parskip}{0pt}%
  \def\l@section##1##2{{\setlength\@tempdima{1.5em}\leftskip1.5em\rightskip\@pnumwidth\parfillskip-\@pnumwidth
    \parindent\z@\noindent{\bfseries ##1}\nobreak
    \leaders\hbox{$\m@th\mkern\@dotsep mu\hbox{.}\mkern\@dotsep mu$}\hfill
    \nobreak\hb@xt@\@pnumwidth{\hss ##2}\par}\vspace{1pt}}%
  \def\l@subsection##1##2{}%
  \def\appendixcontents{\vspace{3pt}{\leftskip1.5em\parindent\z@\noindent\bfseries Appendices\par}\vspace{1pt}}%
  \@starttoc{toc}\endgroup}
\makeatother

\pagestyle{plain}
\setlength{\parindent}{0pt}
\setlength{\parskip}{5pt}
% Final-report figures (both data fractions; scripts/final_figures.py) first, then the mid-project ones that did
% not change (scripts/report_figures.py: the physical system, the pipeline check, the data loader).
\graphicspath{{figures/final/}{figures/}}
% Pages holding only floats fill from the top instead of spreading the floats over the page.
\makeatletter
\setlength{\@fptop}{0pt}
\setlength{\@fpsep}{14pt plus 2pt}
\setlength{\@fpbot}{0pt plus 1fil}
\makeatother
% Let figures and tables share pages with text instead of collecting on half-empty float pages.
\renewcommand{\topfraction}{0.85}
\renewcommand{\bottomfraction}{0.6}
\renewcommand{\textfraction}{0.12}
\renewcommand{\floatpagefraction}{0.8}

% Include a figure if it exists, otherwise a placeholder box naming it.
\newcommand{\figureorpending}[2]{\IfFileExists{figures/final/#1}{\includegraphics[#2]{#1}}{%
    \IfFileExists{figures/#1}{\includegraphics[#2]{#1}}{%
    \fbox{\parbox{0.9\linewidth}{\centering\vspace{1.5em}\textcolor{gray}{\textit{pending: \texttt{\detokenize{#1}}}}\vspace{1.5em}}}}}}

\begin{document}
\pagenumbering{gobble} % the title page is unnumbered; numbering starts at 1 on the next page

\begin{center}
    {\large EECS 221: Machine Learning for Science\par}
    \vspace{5pt}
    {\LARGE\bfseries
    Data--Compute Trade-offs\\[3pt]
    in Neural PDE Forecasting\par}
    \vspace{7pt}
    {\Large Suha Kagan Kose\par}
    \vspace{6pt}
    {\large Final Report\par}
\end{center}

\textbf{Project Goal.} This project compares how accurately and efficiently three neural-network families
(CNextU-Net, TFNO and AViT) forecast Rayleigh--B\'enard convection when training data and compute are limited: 10\% or
100\% of the training trajectories, and one or four hours on one A100 GPU. Each model is also measured against its cost,
a numerical solver, noisy inputs, and how far its forecasts can be trusted.

\textbf{Main findings.}
\begin{enumerate}\setlength{\itemsep}{0pt}\setlength{\parskip}{1pt}
    \item \textbf{The pipeline is validated:} released checkpoints reproduce the paper's test VRMSE (0.2\%, 1.0\%).
    \item \textbf{CNextU-Net is the best after four hours} at both data fractions, one step ahead and over rollouts.
    \item \textbf{Ten times more data helps little} ($-3$\% to $-16.5$\%); four hours instead of one helps more.
    \item \textbf{The ranking depends on the measure:} persistence wins in onset, TFNO on heat transport.
    \item \textbf{A coarse solver matches CNextU-Net at $\mathrm{Ra}\le10^8$;} training pays off after 6,200+ forecasts.
    \item \textbf{Conformal intervals scaled by an ensemble's spread} stay calibrated through 30-step rollouts.
    \item \textbf{Cost:} 146.4 GPU-hours (5,111 SU); the 12 core runs held 83, below the 100 estimated.
\end{enumerate}

\textbf{Structure of the report.} The rest of the report is organised as follows.\par\vspace{4pt}
\compactcontents

\newpage
\pagenumbering{arabic}
\section{Data and Methods}
\label{sec:setup}

\textbf{Physical system} (Fig.~\ref{fig:system}). Rayleigh--B\'enard convection is a horizontal layer of fluid
heated from below and cooled from above. Warm fluid near the bottom is lighter than the cold fluid above it; once
this buoyancy forcing is strong enough to overcome the damping by viscosity and heat diffusion, the layer overturns
into rising warm and sinking cold plumes, with thin thermal boundary layers along both walls. It is the classical
model of buoyancy-driven convection, as in the atmosphere, the oceans, stellar interiors and the Earth's mantle,
and a standard test problem for turbulence and heat transfer. We consider its two-dimensional form: buoyancy $b$
(the temperature variable), velocity $(u, w)$ and pressure $p$ in a $4\times1$ box, periodic in $x$, between a hot
bottom wall ($b=1$) and a cold top wall ($b=0$). Two dimensionless numbers set a regime: the \emph{Rayleigh number}
$\mathrm{Ra}$, the strength of the buoyancy forcing relative to viscous and thermal damping, and the
\emph{Prandtl number} $\mathrm{Pr}$, how fast momentum diffuses relative to heat (equations in
Appendix~\ref{app:rb_regimes}). A larger Rayleigh number means stronger, more turbulent convection with finer plumes
and thinner boundary layers. Started from a fluid at rest, convection first has to grow out of the nearly still state
(\emph{onset}) and then settles into statistically steady, \emph{developed} convection.

\begin{figure}[!htbp]
\centering
\figureorpending{fig_system.pdf}{width=\linewidth}
\caption{The physical system: true buoyancy (temperature) in two trajectories of The Well's test split, both at
$\mathrm{Pr}=1$, at frame 83 of 200 (developed convection). The bottom wall is hot ($b=1$), the top wall is cold
($b=0$), gravity points down and $x$ is periodic. Hot plumes rise from the bottom wall and cold plumes sink from
the top; at the higher Rayleigh number the plumes break up into much finer structures.}
\label{fig:system}
\end{figure}

\textbf{Data.} \texttt{rayleigh\_benard} from The Well~\cite{ohana2024well,wellrb}: 1,750 trajectories, 50 in each
of 35 regimes (five Rayleigh numbers from $10^6$ to $10^{10}$ $\times$ seven Prandtl numbers from 0.1 to 10;
Appendix~\ref{app:rb_regimes}); 367\,GB in all. Each trajectory starts from rest, with a linear buoyancy profile plus
a small random perturbation; within a regime, trajectories differ in the profile's amplitude and the perturbation's
random seed. Each has 200 frames, 0.25 simulation time units apart, of the four fields on a $512\times128$ grid:
uniform and periodic in $x$, and on Chebyshev--Gauss nodes in $y$. We keep The Well's official split, which assigns
whole trajectories: 1,400 for training, 175 for validation and 175 for testing (40 / 5 / 5 in each regime). The
study trains on two fractions of the training split: 10\% (140 trajectories, 4 of the 40 in each regime, 29\,GB)
and 100\% (all 1,400, 294\,GB); the 10\% subset is a fixed draw shared by every model and seed. Every model is
evaluated on the full test split. Fields are normalised with The Well's training statistics.

\textbf{Task.} The study asks which of three neural-network families (Table~\ref{tab:models}) forecasts
Rayleigh--B\'enard convection most usefully when training data and compute are limited. Each model learns the same
forecasting task: from four consecutive frames of buoyancy, pressure and both velocity components it predicts the
next frame, and longer forecasts come from feeding its predictions back in (autoregressive rollout); each trajectory
gives 196 one-step windows.

\textbf{Models.} The Well's own implementations, trained from random initialisation, with matched parameter
counts (Table~\ref{tab:models}). All three use GELU activations. The paper-era code skipped TFNO's Fourier layers, an
error fixed upstream in November 2025; we scored the released TFNO weights with the forward pass they were trained
with, and our own TFNO, trained with the fixed code, uses its Fourier layers (Appendix~\ref{app:tfno}).

\begin{table}[!htbp]
\centering\footnotesize
\begin{tabular*}{\linewidth}{@{\extracolsep{\fill}}llrl@{}}
\toprule
Model & Family & Parameters & Relation to released checkpoint \\
\midrule
CNextU-Net~\cite{liu2022convnext} & ConvNeXt U-Net & 18.58M & same architecture \\
TFNO~\cite{li2021fno,kossaifi2024tfno} & Fourier operator & 19.28M & same configuration, fixed forward \\
AViT~\cite{mccabe2024mpp} & axial-attention ViT & 18.57M & none released; width 384 $\to$ 432 \\
\bottomrule
\end{tabular*}
\caption{The three model families compared in this study.}
\label{tab:models}
\end{table}

\textbf{Training.} Every model is trained the same way: AdamW on the one-step MSE of the normalised fields, batch
64, fp32, on one A100. The learning rate warms up to $10^{-3}$ and then decays to zero along a cosine over four hours
of active GPU training. The core study is $3\times2\times2=12$ runs: three models, two data fractions (10\% and
100\%) and two seeds (initial weights and data order). Each run is evaluated after one and four hours, the two
budgets of the study's compute axis; the one-hour result is the same run's checkpoint, when the learning rate is
still about 85\% of its peak (Appendix~\ref{app:recipe}).

\textbf{References and baselines.} The method is a controlled comparison: the three families are trained from
scratch with one recipe on the same data and budgets, and compared on accuracy and on cost to train and to use.
Three references are shown next to them. (a) \emph{Persistence}, the baseline to beat, predicts that the next frame
equals the last input frame. It needs no training and almost no computation, so a model is worth its training cost
only where it beats persistence; it is hard to beat where the flow changes slowly, as during onset. (b) The Well's
released \emph{CNextU-Net}, trained on all the training data for 12 hours on an H100, shows what this architecture
reaches with a longer schedule; it is a reference point, not a baseline to beat, and it also validated our pipeline
(Section~\ref{sec:m1}). The released TFNO is not used as a reference, since it was trained without its Fourier
layers (Appendix~\ref{app:tfno}). (c) A \emph{numerical solver} (Dedalus~\cite{burns2020dedalus}) restarted from the
last input frame on coarser grids, as a check against weak baselines~\cite{mcgreivy2024weak}
(Section~\ref{sec:solver}).

\textbf{Evaluation and VRMSE.} For one predicted frame and one field $f$, The Well's metric code computes
\[
\mathrm{VRMSE}_f=\sqrt{\frac{\frac{1}{N}\sum_{i=1}^{N}(\hat y_{f,i}-y_{f,i})^2}{s_f^2+\varepsilon}},
\qquad s_f^2=\frac{1}{N-1}\sum_{i=1}^{N}(y_{f,i}-\bar y_f)^2,\qquad \varepsilon=10^{-7},
\]
over the $N=512\times128$ grid points. It compares the size of the error with how much the true frame varies across
the grid; $\varepsilon$ only keeps it from dividing by zero. Scores are averaged over the four fields (buoyancy $b$,
pressure $p$, velocities $u$, $w$) and over windows (one-step: all 196 windows of each of the 175 test trajectories;
rollouts: at each step). The main metric is VRMSE on the fields as stored (\emph{physical-field VRMSE}). The paper
first rescales each field by its training standard deviation (\emph{paper protocol}); during onset this makes
$\varepsilon$ negligible and the scores much larger, so we use the paper protocol only where we compare with the
paper (Appendix~\ref{app:metrics}).

\textbf{Onset and developed convection.} During onset the fluid is nearly still, velocity barely varies across
the grid, and any error is large relative to it; these windows dominate any all-window mean. We therefore report the
two phases separately, with a cut taken from each trajectory's own data: its onset frame $t_{\mathrm{on}}$ is the first
frame at which the domain-mean kinetic energy reaches 10\% of its mean over frames 150--199. A one-step window is
\emph{onset} if its first input frame precedes $t_{\mathrm{on}}$, and \emph{developed} otherwise. On the test split
$t_{\mathrm{on}}$ lies between frames 41 and 78 (median 52), later at low Rayleigh number (median 60 at
$\mathrm{Ra}=10^6$, 48--49 from $10^8$ on), and 26.7\% of the windows are onset. The fixed cut at frame 50 of the
mid-project report is kept as a check: it changes the onset means by 4--5\% and the developed means by at most 3\%,
and no ranking. Rollouts start at frames 20 (onset for every trajectory), 80 and 140 (developed for every trajectory).
Table~\ref{tab:headline} shows why a single all-window number would mislead: three comparisons of this study reverse
between all windows and developed convection, at both data fractions.

\begin{table}[!htbp]
\centering\footnotesize
\begin{tabularx}{\linewidth}{@{}>{\raggedright\arraybackslash}p{0.25\linewidth}l>{\raggedright\arraybackslash}X>{\raggedright\arraybackslash}X@{}}
\toprule
Comparison (4 h) & Data & All windows & Developed convection \\
\midrule
AViT against CNextU-Net & 10\% & 0.134 against 0.136: level & 0.0761 against 0.0505: CNextU-Net 1.5 times better \\
                        & 100\% & 0.112 against 0.130: AViT looks better & 0.0636 against 0.0489: CNextU-Net 1.3 times better \\
\addlinespace[2pt]
Our CNextU-Net against the released one & 10\% & 0.136 against 0.224: ours looks better & 0.0505 against 0.0382:
the released one 1.3 times better \\
                        & 100\% & 0.130 against 0.224: ours looks better & 0.0489 against 0.0382: the released one 1.3
times better \\
\addlinespace[2pt]
Persistence against TFNO & 10\% & 0.199 against 0.209: persistence looks better & 0.229 against 0.111: TFNO 2.1 times
better \\
                         & 100\% & 0.199 against 0.199: level & 0.229 against 0.107: TFNO 2.1 times better \\
\bottomrule
\end{tabularx}
\caption{How the all-window mean can reverse comparisons: physical-field one-step test VRMSE (Table~\ref{tab:test}),
onset and developed convection split at each trajectory's onset frame. The paper protocol, whose all-window mean is
even more dominated by onset, ranks AViT ahead of CNextU-Net at both fractions (Appendix~\ref{app:metrics}).}
\label{tab:headline}
\end{table}

\textbf{Additional measures.} (a) Failure horizon: the first rollout step at which the field-mean VRMSE exceeds 1, 2
or 5. A field held constant at the target's spatial mean scores VRMSE $\approx 1$, so above threshold 1 a forecast
says less about the spatial pattern than that mean alone. We report the fraction of the 525 rollouts (175
trajectories $\times$ 3 starts) that stay below each threshold through step 30, and for threshold 1 its value at every
step (Kaplan--Meier survival). (b) The kinetic-energy spectrum along $x$, to see at which scales the error lies.
(c) The heat transport (Nusselt number), averaged over each 30-step rollout and compared with the truth per Rayleigh
number. (d) Inference time and memory. Input-noise robustness and uncertainty have their own sections
(\ref{sec:noise} and \ref{sec:uq}); definitions are in Appendix~\ref{app:metrics}.

\section{Pipeline Validated Against the Released Checkpoints}
\label{sec:m1}

We loaded The Well's released CNextU-Net and TFNO weights and scored them with our own data path and metric on
every window of the 175 test trajectories (Table~\ref{tab:released}). Under the paper protocol both reproduce the
published one-step values, and their rollouts from frame 0 are above 10, as published. The released TFNO reproduces
its value only with the forward pass it was trained with, which skipped its Fourier layers: it is effectively a
per-pixel network that cannot move structures across the grid, and its error tracks persistence
(Fig.~\ref{fig:released}). The figure also shows why the phases must be reported separately: under the paper
protocol, onset windows give 96\% of the released CNextU-Net's all-window mean (Appendix~\ref{app:metrics}).

\begin{table}[!htbp]
\centering\small
\begin{tabular}{@{}lcccc@{}}
\toprule
& \multicolumn{2}{c}{Paper-protocol one-step VRMSE} & \multicolumn{2}{c}{Rollout from frame 0, steps 6--12 / 13--30} \\
\cmidrule(lr){2-3}\cmidrule(l){4-5}
Model & Ours (released weights) & Published & Ours & Published \\
\midrule
CNextU-Net & 0.6688 & 0.6699 & 36 / 35 & $>10$ / $>10$ \\
TFNO$^{*}$ & 0.6634 & 0.6566 & 142 / 169 & $>10$ / $>10$ \\
\bottomrule
\end{tabular}
\caption{The released CNextU-Net and TFNO scored by our pipeline on the test split, against the published results.
$^{*}$With the paper-era forward pass that skipped its Fourier layers (Appendix~\ref{app:tfno}).}
\label{tab:released}
\end{table}

\begin{figure}[!htbp]
\centering
\figureorpending{fig_released_by_window.pdf}{scale=1}
\caption{Released checkpoints' one-step VRMSE (paper protocol) on the test set, by each window's first input frame
(mean over the 175 trajectories), against persistence. Shaded: the first 50 frames, where most trajectories are
still in onset.}
\label{fig:released}
\end{figure}

\section{Training Setup and Budgets}
\label{sec:m2}

\textbf{Learning rate.} The paper's rates ($5\cdot10^{-4}$ for CNextU-Net, $10^{-4}$ for TFNO) were chosen for
12-hour runs. With one to four hours a model takes far fewer optimizer steps, so we tested whether a higher rate
suits our budgets. An equal-budget search (20 minutes of GPU training per run, 10\% data, validation only) chose
$10^{-3}$ for all three models (Appendix~\ref{app:recipe}): 0.408 against 0.469 for CNextU-Net at
$5\cdot10^{-4}$, and 0.302 against 0.603 for TFNO at $10^{-4}$. It is the largest rate tested.

\textbf{Batch size and precision.} Batch 64 (The Well's) fits every model in fp32 on an A100 80\,GB; CNextU-Net
peaks at 64.8\,GB of memory reserved by PyTorch (60.9\,GB allocated, Table~\ref{tab:resources}). TFNO's Fourier
layers fail in bf16, so all models use fp32. This favours TFNO: bf16 would speed up a training step 1.7 times for
CNextU-Net and 2.2 times for AViT.

\textbf{Training budget.} Both budgets count only \emph{active GPU training time}: the clock runs while the GPU
computes training steps (forward pass, backward pass and weight update) and stops otherwise. Elapsed time would partly
measure the cluster's shared storage, which cannot feed TFNO and AViT fast enough (Appendix~\ref{app:accounting}).
So every run takes the same number of optimizer steps at either data fraction (about 12,100 for CNextU-Net, 26,300
for TFNO and 74,900 for AViT in four hours), and \emph{allocated GPU time}, the time the job held the GPU, is
reported as a cost. The 10\% runs read their data from the shared file system; the 100\% runs first copied the
training split to the GPU node's own disk (4--8 minutes), which removed almost all of their data waits.

\section{Accuracy Under Limited Data and Compute}
\label{sec:results}

This section answers the project's question for accuracy: how accurate each family becomes with 10\% and 100\% of
the training data after one and four hours of training. Section~\ref{sec:cost} sets these results against cost.

\textbf{Training} (Fig.~\ref{fig:curves}). All twelve runs completed their 4-hour budgets. On a fixed validation
sample of 64 windows, AViT and CNextU-Net reach similar errors after four hours (0.14--0.17), TFNO a higher one
(0.27--0.32), and TFNO's loss levels off after about 1.5 hours at both fractions. With all the data each model makes
a tenth as many passes over the training windows (CNextU-Net 2.8 against 28, TFNO 6.1 against 61, AViT 17.5 against
175), yet its training loss follows almost the same curve.

\begin{figure}[!htbp]
\centering
\figureorpending{fig_training_f010.pdf}{scale=1}\\[2pt]
\figureorpending{fig_training_f100.pdf}{scale=1}
\caption{Training MSE in normalised units (left) and validation VRMSE on the 64 fixed windows (right), against
active training time: 10\% of the training data (top) and 100\% (bottom). Solid: seed 0; dashed: seed 1. Vertical
lines: the 1-hour and 4-hour checkpoints.}
\label{fig:curves}
\end{figure}

\textbf{One-step test error} (all 175 test trajectories, both seeds; Table~\ref{tab:test}, Fig.~\ref{fig:testwindow}).
After four hours the two phases separate the models. In developed convection CNextU-Net is the most accurate at both
fractions (0.0505 and 0.0489), ahead of AViT (0.0761 and 0.0636) and TFNO (0.1105 and 0.1074); during onset AViT is
the most accurate. The seed-to-seed standard deviation is at most 0.006 over all windows and 0.0004 in developed
convection. The released CNextU-Net remains better in developed convection (0.0382), but ours is better during onset
(0.353 against 0.733). Persistence is better than every model during onset (0.116) and 2--4.7 times worse in
developed convection (0.229). After one hour the picture differs: AViT has the lowest all-window error at 10\% data
and ties CNextU-Net at 100\%, and CNextU-Net varies most between seeds.

\begin{table}[!htbp]
\centering\small
\begin{tabular}{@{}llcccc@{}}
\toprule
& & \multicolumn{3}{c}{Physical-field one-step VRMSE} & Paper protocol \\
\cmidrule(lr){3-5}
Model, checkpoint & Data & all windows & onset & developed & (all windows) \\
\midrule
CNextU-Net, 4 h & 10\% & 0.136 $\pm$ 0.006 & 0.370 $\pm$ 0.024 & \textbf{0.0505} $\pm$ 0.0002 & 0.349 $\pm$ 0.022 \\
                & 100\% & 0.130 $\pm$ 0.001 & 0.353 $\pm$ 0.003 & \textbf{0.0489} $\pm$ 0.0004 & 0.333 $\pm$ 0.002 \\
TFNO, 4 h & 10\% & 0.209 $\pm$ 0.001 & 0.480 $\pm$ 0.002 & 0.1105 $\pm$ 0.0001 & 0.488 $\pm$ 0.001 \\
          & 100\% & 0.199 $\pm$ 0.003 & 0.452 $\pm$ 0.014 & 0.1074 $\pm$ 0.0003 & 0.459 $\pm$ 0.015 \\
AViT, 4 h & 10\% & \textbf{0.134} $\pm$ 0.001 & \textbf{0.294} $\pm$ 0.003 & 0.0761 $\pm$ 0.0002 & \textbf{0.298} $\pm$ 0.002 \\
          & 100\% & \textbf{0.112} $\pm$ 0.000 & \textbf{0.245} $\pm$ 0.000 & 0.0636 $\pm$ 0.0001 & \textbf{0.245} $\pm$ 0.000 \\
\midrule
CNextU-Net, 1 h & 10\% & 0.728 $\pm$ 0.176 & 2.510 $\pm$ 0.653 & \textbf{0.0783} $\pm$ 0.0022 & 2.345 $\pm$ 0.599 \\
                & 100\% & \textbf{0.456} $\pm$ 0.035 & 1.503 $\pm$ 0.131 & \textbf{0.0751} $\pm$ 0.0001 & 1.393 $\pm$ 0.106 \\
TFNO, 1 h & 10\% & 0.473 $\pm$ 0.078 & 1.442 $\pm$ 0.292 & 0.1203 $\pm$ 0.0001 & 1.362 $\pm$ 0.249 \\
          & 100\% & 0.491 $\pm$ 0.047 & 1.494 $\pm$ 0.183 & 0.1262 $\pm$ 0.0023 & 1.417 $\pm$ 0.161 \\
AViT, 1 h & 10\% & \textbf{0.394} $\pm$ 0.036 & \textbf{1.244} $\pm$ 0.135 & 0.0846 $\pm$ 0.0001 & \textbf{1.171} $\pm$ 0.125 \\
          & 100\% & \textbf{0.456} $\pm$ 0.062 & \textbf{1.485} $\pm$ 0.234 & 0.0804 $\pm$ 0.0001 & \textbf{1.389} $\pm$ 0.209 \\
\midrule
Released CNextU-Net & 100\% & 0.224 & 0.733 & 0.0382 & 0.669 \\
Persistence & -- & 0.199 & 0.116 & 0.2292 & 0.202 \\
\bottomrule
\end{tabular}
\caption{One-step test error: mean $\pm$ standard deviation over the two seeds (0.000: below 0.0005). Onset and
developed split at each trajectory's onset frame. Bold: best of our models at each budget and data fraction.}
\label{tab:test}
\end{table}

\begin{figure}[!htbp]
\centering
\figureorpending{fig_test_by_window.pdf}{scale=1}
\caption{One-step test VRMSE by window, mean over the 175 test trajectories: our models after four hours (seed mean,
band: seed range) trained on (a) 10\% and (b) 100\% of the data, the released CNextU-Net and persistence. Dark
shading: before the earliest onset frame; light shading: between the earliest and the latest, where the phase
depends on the trajectory.}
\label{fig:testwindow}
\end{figure}

\textbf{Error around each trajectory's onset} (Fig.~\ref{fig:onset}). Counted from each trajectory's own onset frame
$t_{\mathrm{on}}$, the error curves of the five Rayleigh numbers line up. Our CNextU-Net's onset error (above twice
its level in developed convection) ends 9--22 frames before $t_{\mathrm{on}}$ at every Rayleigh number, against frames
26--59 counted from frame 0; for the released CNextU-Net, 6--19 frames before $t_{\mathrm{on}}$ against frames 27--63.
The large errors therefore lie in the first part of the onset windows: every model beats persistence from 17--23
frames before $t_{\mathrm{on}}$, so persistence wins the onset mean on the earliest, nearly still frames alone.
Just after $t_{\mathrm{on}}$, while convection grows, the error rises a second time and peaks 9--16 frames later
(CNextU-Net and TFNO by 32--40\%, AViT by 70--90\%, the released model by 36\%); persistence's error peaks there too
(14 frames after $t_{\mathrm{on}}$), because the flow changes fastest. The residual trend in Fig.~\ref{fig:onset}c,
an earlier drop at higher Rayleigh number, means that the kinetic-energy fraction at which the error ends depends
somewhat on the regime: at $\mathrm{Ra}=10^{10}$ about 20 low-error windows per trajectory count as onset, at $10^6$
about 10.

\begin{figure}[!htbp]
\centering
\figureorpending{fig_onset_aligned.pdf}{scale=1}
\caption{One-step test VRMSE around convection onset, four-hour models. (a) Against the window's start frame minus
its trajectory's onset frame $t_{\mathrm{on}}$, mean over the 175 test trajectories (seed mean; solid: 100\% data,
dashed: 10\%); shown where every trajectory contributes. (b, c) Our 100\% CNextU-Net per Rayleigh number (35
trajectories each), (b) from frame 0 and (c) from $t_{\mathrm{on}}$.}
\label{fig:onset}
\end{figure}

\textbf{Ten times more data} (Fig.~\ref{fig:scaling}). After four hours, going from 10\% to 100\% of the data lowers
the developed-convection error by 3.1\% for CNextU-Net, 2.8\% for TFNO and 16.5\% for AViT; AViT improves on every
field (buoyancy $-10$\%, velocity $-19$ to $-21$\%, pressure $-21$\%) and during onset ($-17$\%). The ranking in
developed convection is unchanged, but AViT's gap to CNextU-Net narrows from 1.51 to 1.30 times. Training time matters
more than data at these budgets: from one to four hours CNextU-Net's developed error falls by 35\% at both fractions,
and with all the data AViT and TFNO gain about twice as much from the extra hours as at 10\% ($-21$\% against $-10$\%,
$-15$\% against $-8$\%). A reading consistent with this, not tested: at 10\% AViT and TFNO have seen their subset many
times and are limited by data, while CNextU-Net, with few passes, is limited by compute at both fractions. The
one-step gain does not reach the rollouts: after 30 steps from developed flow no model is better with more data
(Fig.~\ref{fig:scaling}c). Ten times more data also does not close our CNextU-Net's gap to the released one in
developed convection (1.32 to 1.28 times); the released model trained for about 12 hours on an H100.

\begin{figure}[!htbp]
\centering
\figureorpending{fig_data_scaling.pdf}{scale=1}
\caption{Data scaling: test VRMSE against the fraction of training data, after one hour (open markers, dashed) and
four hours (filled, solid), seed mean with the seed range. (a) One step, developed convection; (b) one step, onset;
(c) after 30 rollout steps from frames 80 and 140. Dashed black: released CNextU-Net; dotted: persistence.}
\label{fig:scaling}
\end{figure}

\textbf{Rollouts} (Table~\ref{tab:rollout}, Fig.~\ref{fig:rollouts}). Our models were trained to predict one step
ahead from true frames; in a rollout they receive their own slightly wrong outputs~\cite{brandstetter2022mp}, so
errors can grow from step to step~\cite{mccabe2023stability} and the best one-step model need not be the best over a
rollout. After four hours CNextU-Net is the most stable at both fractions: after 30 steps from developed flow its
error is 0.52--0.53 and 0.40 (from frames 80 and 140), below persistence's 0.86 and 0.54, and 88--89\% of its rollouts
stay below VRMSE 1 throughout. It is also more stable than the released CNextU-Net (0.63 and 0.51; 77\%), although
the released model is better one step ahead. AViT, nearly as accurate as CNextU-Net one step ahead, degrades fastest:
its mean error crosses persistence's after 21--25 steps, and only 61--63\% of its rollouts stay below VRMSE 1. TFNO
lies in between. After one hour the ranking depends on the data: at 10\% TFNO is the most stable (60\% of rollouts
below VRMSE 1, against 34\%), at 100\% TFNO is the least stable (23\%, in both seeds) and CNextU-Net and AViT reach
50--52\%. We have no explanation for TFNO's unstable one-hour rollouts at 100\%; after four hours they match the 10\%
runs. From frame 20, in onset, persistence beats every model.

\begin{table}[!htbp]
\centering\footnotesize
\setlength{\tabcolsep}{5pt}
\begin{tabular}{@{}llccc@{}}
\toprule
& & VRMSE after 30 steps & Fraction staying below threshold & Paper protocol, steps \\
Model, checkpoint & Data & from frame 20 / 80 / 140 & through step 30: 1 / 2 / 5 & 6--12 / 13--30 (normalised) \\
\midrule
CNextU-Net, 1 h & 10\% & 3.60 / 1.09 / 1.04 & 0.34 / 0.74 / 0.79 & 202 / 148 \\
                & 100\% & 2.78 / 1.00 / 0.86 & 0.50 / 0.76 / 0.82 & 114 / 103 \\
TFNO, 1 h & 10\% & 2.67 / 0.90 / 0.78 & 0.60 / 0.77 / 0.86 & 81 / 116 \\
          & 100\% & 3.58 / 1.23 / 1.09 & 0.23 / 0.76 / 0.83 & 100 / 150 \\
AViT, 1 h & 10\% & 4.67 / 1.11 / 0.98 & 0.34 / 0.74 / 0.83 & 129 / 184 \\
          & 100\% & 5.26 / 0.97 / 0.85 & 0.52 / 0.75 / 0.82 & 157 / 217 \\
\midrule
CNextU-Net, 4 h & 10\% & 1.78 / \textbf{0.53} / \textbf{0.40} & \textbf{0.89} / \textbf{0.94} / \textbf{0.98} & \textbf{15} / \textbf{30} \\
                & 100\% & 1.78 / \textbf{0.52} / \textbf{0.40} & \textbf{0.88} / \textbf{0.92} / \textbf{0.98} & \textbf{19} / \textbf{34} \\
TFNO, 4 h & 10\% & 1.79 / 0.59 / 0.45 & 0.75 / 0.85 / 0.91 & 51 / 77 \\
          & 100\% & 1.89 / 0.57 / 0.44 & 0.77 / 0.86 / 0.92 & 44 / 63 \\
AViT, 4 h & 10\% & 5.42 / 0.80 / 0.65 & 0.63 / 0.80 / 0.87 & 75 / 167 \\
          & 100\% & 5.64 / 0.83 / 0.69 & 0.61 / 0.81 / 0.89 & 61 / 148 \\
\midrule
Released CNextU-Net & 100\% & 1.65 / 0.63 / 0.51 & 0.77 / 0.82 / 0.90 & 36 / 35 \\
Persistence & -- & 0.68 / 0.86 / 0.54 & 0.94 / 1.00 / 1.00 & 0.27 / 0.50 \\
\bottomrule
\end{tabular}
\caption{Rollouts on the test set, physical-field VRMSE unless stated, mean over the two seeds. Middle column:
fraction of the 525 rollouts (175 trajectories $\times$ 3 start frames) whose error stays below each threshold at
every step through step 30. Right: the paper's Table 3 protocol, rollouts from frame 0 (published: $>10$ in both
windows). Bold: best of our models at four hours.}
\label{tab:rollout}
\end{table}

\begin{figure}[!htbp]
\centering
\figureorpending{fig_rollouts.pdf}{scale=1}
\caption{Rollouts of our models (seed mean; solid: 100\% data, dashed: 10\%), the released CNextU-Net and
persistence, after one hour (a, b) and four hours (c, d). (a, c) VRMSE at each step, mean over the rollouts from
frames 80 and 140 (developed convection). (b, d) Fraction of all 525 rollouts that have stayed below VRMSE 1 up to
each step (Kaplan--Meier survival).}
\label{fig:rollouts}
\end{figure}

\textbf{Field-wise error} (Table~\ref{tab:fields}). In developed convection buoyancy is the hardest field for every
model (CNextU-Net: 0.095, against 0.022--0.025 for pressure and about 0.04 for each velocity component), consistent
with the sharp gradients of its thin plumes and boundary layers; CNextU-Net's buoyancy error does not change with
data, while its velocity and pressure errors improve by 4--11\%. During onset the error comes almost entirely from
velocity, whose near-zero spatial variance magnifies relative error; on buoyancy all three models beat persistence
even then (0.026--0.039 against 0.058), so their onset deficit against persistence is a velocity effect. After 30
steps from developed flow, CNextU-Net and TFNO stay below persistence on every field, while AViT is worse than
persistence on horizontal velocity (1.06--1.12 against 0.93).

\begin{table}[!htbp]
\centering\footnotesize
\setlength{\tabcolsep}{3.0pt}
\begin{tabular}{@{}llcccccccccccc@{}}
\toprule
& & \multicolumn{4}{c}{One step, developed} & \multicolumn{4}{c}{One step, onset} & \multicolumn{4}{c}{30 steps from frames 80, 140} \\
\cmidrule(lr){3-6}\cmidrule(lr){7-10}\cmidrule(l){11-14}
Model & Data & $b$ & $p$ & $u$ & $w$ & $b$ & $p$ & $u$ & $w$ & $b$ & $p$ & $u$ & $w$ \\
\midrule
CNextU-Net, 4 h & 10\% & \textbf{0.095} & \textbf{0.025} & \textbf{0.039} & \textbf{0.044} & 0.028 & 0.011 & 0.827 & 0.612 & \textbf{0.384} & \textbf{0.293} & \textbf{0.598} & \textbf{0.584} \\
                & 100\% & \textbf{0.095} & \textbf{0.022} & \textbf{0.037} & \textbf{0.042} & 0.029 & 0.011 & 0.786 & 0.588 & \textbf{0.379} & \textbf{0.288} & \textbf{0.589} & \textbf{0.572} \\
TFNO, 4 h & 10\% & 0.174 & 0.055 & 0.096 & 0.117 & 0.039 & 0.015 & 1.117 & 0.750 & 0.466 & 0.307 & 0.665 & 0.646 \\
          & 100\% & 0.171 & 0.052 & 0.093 & 0.114 & 0.038 & 0.015 & 0.993 & 0.761 & 0.456 & 0.297 & 0.635 & 0.622 \\
AViT, 4 h & 10\% & 0.122 & 0.040 & 0.061 & 0.081 & \textbf{0.027} & \textbf{0.007} & \textbf{0.623} & \textbf{0.518} & 0.497 & 0.419 & 1.059 & 0.922 \\
          & 100\% & 0.109 & 0.031 & 0.050 & 0.064 & \textbf{0.026} & \textbf{0.007} & \textbf{0.519} & \textbf{0.430} & 0.507 & 0.440 & 1.117 & 0.984 \\
\midrule
CNextU-Net, 1 h & 10\% & \textbf{0.129} & \textbf{0.046} & \textbf{0.069} & \textbf{0.069} & 0.046 & 0.033 & 6.553 & 3.409 & 1.017 & 0.729 & 1.280 & 1.249 \\
                & 100\% & \textbf{0.127} & \textbf{0.043} & \textbf{0.063} & \textbf{0.068} & 0.045 & 0.033 & 4.041 & \textbf{1.892} & 0.762 & 0.623 & \textbf{1.211} & 1.117 \\
TFNO, 1 h & 10\% & 0.184 & 0.063 & 0.105 & 0.129 & 0.051 & 0.028 & 3.592 & 2.095 & 0.925 & \textbf{0.603} & \textbf{0.930} & \textbf{0.908} \\
          & 100\% & 0.194 & 0.071 & 0.110 & 0.130 & 0.055 & 0.033 & 3.518 & 2.368 & 1.301 & 0.846 & 1.320 & 1.171 \\
AViT, 1 h & 10\% & 0.132 & 0.047 & 0.071 & 0.088 & \textbf{0.035} & \textbf{0.017} & \textbf{2.965} & \textbf{1.960} & \textbf{0.837} & 0.652 & 1.484 & 1.209 \\
          & 100\% & 0.128 & 0.043 & 0.068 & 0.083 & \textbf{0.036} & \textbf{0.018} & \textbf{3.327} & 2.561 & \textbf{0.698} & \textbf{0.566} & 1.275 & \textbf{1.096} \\
\midrule
Released CNextU-Net & 100\% & 0.071 & 0.020 & 0.029 & 0.033 & 0.021 & 0.011 & 1.404 & 1.497 & 0.505 & 0.366 & 0.713 & 0.697 \\
Persistence & -- & 0.278 & 0.124 & 0.248 & 0.267 & 0.058 & 0.014 & 0.197 & 0.195 & 0.531 & 0.408 & 0.934 & 0.916 \\
\bottomrule
\end{tabular}
\caption{Field-wise test VRMSE (physical fields, mean over the two seeds): buoyancy $b$, pressure $p$, horizontal
and vertical velocity $u$, $w$. Bold: best of our models at each budget and data fraction.}
\label{tab:fields}
\end{table}

\textbf{Physical fidelity} (Fig.~\ref{fig:spectra}). The kinetic-energy spectrum shows how much energy a prediction
has at each scale, from large plumes (small wavenumber $k_x$) to fine detail (large $k_x$); we report predicted over
true energy per wavenumber (1 is exact). All three models put too much energy at the finest scales ($k_x>100$), as is
typical of networks trained with MSE, which barely constrains the weak small
scales~\cite{rahaman2019spectral,lippe2023pderefiner}. At larger scales they differ. CNextU-Net is close to the true
spectrum. TFNO is too smooth: it loses energy from $k_x\approx25$ after one step and from $k_x\approx10$ after 30
steps, down to a tenth of the true energy near $k_x=80$, consistent with its Fourier layers keeping only the coarsest
modes ($|k_x|\le12.6$)~\cite{li2021fno}. AViT is accurate after one step but drifts over the rollout. These spectra do
not change with data: the log spectral error after 30 steps is 1.32 / 1.30 / 2.40 at 10\% and 1.32 / 1.29 / 2.44 at
100\% for CNextU-Net / TFNO / AViT.

\begin{figure}[!htbp]
\centering
\figureorpending{fig_spectra.pdf}{scale=1}
\caption{Horizontal kinetic-energy spectra of our models trained on all the data (four hours). (a) True and
predicted spectrum after 30 rollout steps for one regime ($\mathrm{Ra}=10^8$, $\mathrm{Pr}=1$). (b, c) Predicted
over true spectrum after 1 and 30 steps, per wavenumber: median over the 35 regimes (band: middle 50\%); 1 is exact.
Shaded: small scales ($k_x > 100$).}
\label{fig:spectra}
\end{figure}

\textbf{What the predictions look like} (Fig.~\ref{fig:snapshots}). After 30 steps at $\mathrm{Ra}=10^8$ every model
keeps the behaviour it had with 10\% of the data. CNextU-Net follows the true plumes, though not their exact
positions. TFNO blurs them and builds a spurious warm layer along the bottom wall. AViT shows rectangular artifacts at
the size of its $16\times16$-pixel patches. At $\mathrm{Ra}=10^{10}$ the true
flow is far finer than any of the three reproduces (Fig.~\ref{fig:snapshots10}).

\begin{figure}[!htbp]
\centering
\figureorpending{fig_snapshots_fractions_buoyancy_step30_rayleigh_benard_Rayleigh_1e8_Prandtl_1__0__s080.pdf}{scale=1}
\caption{Buoyancy on a test trajectory at $\mathrm{Ra}=10^8$, $\mathrm{Pr}=1$, after 30 rollout steps from frame 80:
the true flow (top) and our four-hour models (seed 0) trained on 10\% (left) and 100\% (right) of the data. Bright:
warm; dark: cold. One colour range for all panels: the 1st to 99th percentile of the true field away from the walls,
so the hot bottom wall saturates.}
\label{fig:snapshots}
\end{figure}

\begin{figure}[!htbp]
\centering
\figureorpending{fig_snapshots_fractions_buoyancy_step30_rayleigh_benard_Rayleigh_1e10_Prandtl_1__0__s080.pdf}{scale=1}
\caption{As Fig.~\ref{fig:snapshots}, at $\mathrm{Ra}=10^{10}$, $\mathrm{Pr}=1$.}
\label{fig:snapshots10}
\end{figure}

\textbf{Error by Rayleigh number} (Fig.~\ref{fig:rayleigh}a, b). The one-step error in developed convection grows
four- to tenfold from $\mathrm{Ra}=10^6$ to $10^{10}$ for every model as the flow becomes more turbulent (CNextU-Net
0.016 to 0.11); persistence's grows too (0.098 to 0.35). The ranking is the same at every Rayleigh number, and more
data helps AViT at every Rayleigh number. The released CNextU-Net matches ours at $10^6$ and $10^7$ and is better only
at higher Rayleigh numbers. After 30 steps from developed flow, CNextU-Net and TFNO stay below persistence at every
Rayleigh number, while AViT is worse than persistence up to $10^8$ at both fractions.

\textbf{Heat transport} (Table~\ref{tab:nusselt}, Fig.~\ref{fig:rayleigh}c). The Nusselt number is the heat carried
across the layer relative to pure conduction (Appendix~\ref{app:metrics}). We average it over each 30-step rollout
from developed flow, predicted and true alike, and report the median relative error per Rayleigh number. After four
hours TFNO comes closest (median error 20--21\%), with a small bias at every Rayleigh number; below $10^{10}$ it is
closer than persistence (9--29\% against 17--32\%), whose error measures how far one snapshot's heat transport lies
from the 30-step mean. AViT carries too little heat at every Rayleigh number ($-17$\% overall, $-38$ to $-43$\% at
$10^{10}$). Our CNextU-Net is nearly unbiased up to $10^9$ but scattered (median error 49\%), and overshoots by
35--41\% at $10^{10}$. Ten times more data barely changes any of them. The released CNextU-Net carries about twice the
true heat (+108\%, rising from +34\% at $10^6$ to +299\% at $10^{10}$): right at first ($-4$\% at step 1), it drifts to
+100\% by step 12 and +201\% by step 30. A low one-step error therefore does not guarantee the right heat transport
over a rollout, and heat transport ranks the models differently from VRMSE. After one hour
(Fig.~\ref{fig:nusselt}a) our CNextU-Net carries far too much heat (+545\% at 10\%, +133\% at 100\%, growing with the
Rayleigh number to +1026\% and +292\% at $10^{10}$). TFNO and AViT have overall biases within $\pm30$\%, but they
scatter widely (median $|$error$|$ 35--91\%) and reach +60\% and $-62$\% at single Rayleigh numbers; in more than a
quarter of the rollouts of AViT (100\%) at $\mathrm{Ra}\ge10^9$ and of TFNO (10\%) at $10^{10}$, the predicted
Nusselt number is negative, heat carried downwards. The true 30-step average itself fluctuates: between the windows
from frames 80 and 140 of the same trajectory it changes by a median of 15\% (middle 50\%: $-14$\% to $+15$\%), so
errors of that size are within the truth's own variation; across a regime's seven Prandtl numbers and five
trajectories the middle 50\% of the true values spans $-12$\% to $+21$\% of the median (Table~\ref{tab:nusselt}).

\begin{figure}[!htbp]
\centering
\figureorpending{fig_by_rayleigh.pdf}{width=\linewidth}
\caption{Error by Rayleigh number, four-hour models (seed mean; solid: 100\% data, dashed: 10\%; each point
averages seven Prandtl numbers, 35 test trajectories). (a) One-step VRMSE in developed convection. (b) VRMSE after 30
steps from frames 80 and 140. (c) Median signed relative error of the Nusselt number averaged over those 30-step
rollouts. Dashed black: released CNextU-Net; dotted: persistence.}
\label{fig:rayleigh}
\end{figure}

\begin{figure}[!htbp]
\centering
\figureorpending{fig_nusselt.pdf}{scale=1}
\caption{Predicted against true heat transport: the Nusselt number averaged over each 30-step rollout from frames 80
and 140, median over the 70 rollouts per Rayleigh number (seed mean; bars: middle 50\% of the rollouts), after (a) one
and (b) four hours. Black: truth, band: its middle 50\%. The released CNextU-Net is the same in both panels;
persistence (not shown) lies on the truth. Bars that leave the axis: the lower quartile is negative.}
\label{fig:nusselt}
\end{figure}

\begin{table}[!htbp]
\centering\footnotesize
\begin{tabular}{@{}llcccccc@{}}
\toprule
& & \multicolumn{5}{c}{Median signed error per Rayleigh number} & All $\mathrm{Ra}$: \\
\cmidrule(lr){3-7}
Model, checkpoint & Data & $10^6$ & $10^7$ & $10^8$ & $10^9$ & $10^{10}$ & signed / $|$error$|$ \\
\midrule
CNextU-Net, 4 h & 10\% & $-3$\% & $-3$\% & $-9$\% & $-15$\% & $+35$\% & $-1$\% / 49\% \\
                & 100\% & $-8$\% & $+6$\% & $-5$\% & $+9$\% & $+41$\% & $+2$\% / 49\% \\
TFNO, 4 h & 10\% & $-2$\% & $+16$\% & $+19$\% & $-3$\% & $-14$\% & $+6$\% / \textbf{20}\% \\
          & 100\% & $-3$\% & $+15$\% & $+17$\% & $+5$\% & $-13$\% & $+8$\% / \textbf{21}\% \\
AViT, 4 h & 10\% & $-11$\% & $-8$\% & $-13$\% & $-18$\% & $-43$\% & $-17$\% / 27\% \\
          & 100\% & $-5$\% & $-17$\% & $-19$\% & $-14$\% & $-38$\% & $-17$\% / 25\% \\
\midrule
CNextU-Net, 1 h & 10\% & $+204$\% & $+363$\% & $+673$\% & $+926$\% & $+1026$\% & $+545$\% / 548\% \\
                & 100\% & $+38$\% & $+83$\% & $+163$\% & $+214$\% & $+292$\% & $+133$\% / 153\% \\
TFNO, 1 h & 10\% & $+9$\% & $+38$\% & $+35$\% & $-5$\% & $-43$\% & $+21$\% / 91\% \\
          & 100\% & $-2$\% & $+28$\% & $+60$\% & $+34$\% & $+43$\% & $+29$\% / 48\% \\
AViT, 1 h & 10\% & $-2$\% & $+18$\% & $+30$\% & $+34$\% & $+20$\% & $+14$\% / \textbf{35}\% \\
          & 100\% & $-14$\% & $-10$\% & $-27$\% & $-55$\% & $-62$\% & $-28$\% / 44\% \\
\midrule
Released CNextU-Net & 100\% & $+34$\% & $+73$\% & $+121$\% & $+165$\% & $+299$\% & $+108$\% / 114\% \\
Persistence & -- & $-2$\% & $-3$\% & $+6$\% & $+7$\% & $+6$\% & $+2$\% / 27\% \\
\midrule
True Nusselt number: median & & 7.9 & 13.3 & 23.2 & 44.1 & 96.1 & \\
\quad middle 50\% & & 7.2--8.4 & 11.9--14.8 & 20.3--26.5 & 40.9--50.7 & 88.2--115.9 & \\
\bottomrule
\end{tabular}
\caption{Relative error of the Nusselt number averaged over each 30-step rollout from frames 80 and 140 (70 rollouts
per Rayleigh number), median over rollouts, mean over the two seeds. Negative: too little heat transport. Bold: closest
of our models at each checkpoint.}
\label{tab:nusselt}
\end{table}

\section{Accuracy Against Cost}
\label{sec:cost}

Because data and compute are the limiting resources, we set each model's accuracy against what it costs to train and
to use (Fig.~\ref{fig:cost}, Table~\ref{tab:resources}).
\begin{itemize}
    \item \emph{Training time} (Fig.~\ref{fig:cost}a). From one to four hours the developed-convection error falls
    by 35\% for CNextU-Net at both fractions, by 10\% (10\%) and 21\% (100\%) for AViT, and by 8\% and 15\% for TFNO.
    \item \emph{Cost to train} (Fig.~\ref{fig:cost}b, Fig.~\ref{fig:time}). Every run trained for the same four
    hours, but the GPU time it held differs. Reading from the shared file system, the 10\% runs held 5.2--6.0 hours
    (CNextU-Net), 6.9--7.1 (TFNO) and 11.1--12.2 (AViT), since AViT consumes 330 training samples per second and the
    storage delivers about 98. With the data copied to the GPU node's disk, the 100\% runs held 4.4--4.6, 5.5--5.8 and
    5.7--6.0 hours, although they read ten times more data. AViT needs a seventh of CNextU-Net's memory and takes six
    times as many optimizer steps.
    \item \emph{Cost to use} (Fig.~\ref{fig:cost}c). All three are cheap to run: a 30-step forecast takes
    133--249\,ms with under 0.5\,GB of GPU memory. TFNO is the fastest for one forecast at a time, and AViT the
    cheapest for many at once (0.89\,ms per sample in batches of 64, against 3.15 for TFNO and 6.74 for CNextU-Net).
\end{itemize}

\begin{figure}[!htbp]
\centering
\figureorpending{fig_accuracy_cost.pdf}{width=\linewidth}
\caption{Accuracy against cost (seed mean; filled: 100\% data, open: 10\%). (a) Developed-convection one-step VRMSE
after one and four hours of active training. (b) The same four-hour error against the GPU time each run held
(range over seeds). (c) VRMSE after 30 steps from developed flow against the inference time of a 30-step rollout
(batch 1, one A100).}
\label{fig:cost}
\end{figure}

\begin{table}[!htbp]
\centering\footnotesize
\setlength{\tabcolsep}{4pt}
\begin{tabular}{@{}lcccccccc@{}}
\toprule
& \multicolumn{5}{c}{Training (4-hour budget, batch 64)} & \multicolumn{3}{c}{Inference} \\
\cmidrule(lr){2-6}\cmidrule(l){7-9}
& Optimizer & Epochs & Peak memory, & \multicolumn{2}{c}{GPU time held per run} & 30 steps, & per sample, & Memory, \\
Model & steps & 10\% / 100\% & alloc.\ (reserved) & 10\% & 100\% & batch 1 & batch 64 & batch 1 \\
\midrule
CNextU-Net & $\approx$12,100 & 28 / 2.8 & 60.9 (64.8)\,GB & 5.2--6.0\,h & 4.4--4.6\,h & 248\,ms & 6.74\,ms & 0.31\,GB \\
TFNO & $\approx$26,300 & 61 / 6.1 & 48.7 (57.1)\,GB & 6.9--7.1\,h & 5.5--5.8\,h & 133--135\,ms & 3.15\,ms & 0.44\,GB \\
AViT & $\approx$75,000 & 175 / 17.5 & 8.8 (8.9)\,GB & 11.1--12.2\,h & 5.7--6.0\,h & 163--168\,ms & 0.89\,ms & 0.23\,GB \\
\bottomrule
\end{tabular}
\caption{Cost to train and to use, on one A100 80\,GB in fp32. Epochs: passes over the training windows. Peak memory:
the larger of the two fractions. GPU time held: allocated time of the training job, including data waits, start-up,
the data copy, validation and checkpoints (Appendix~\ref{app:accounting}). Per sample, batch 64: forward pass in the
training profile (with gradient tracking, an upper bound for inference). Inference is the same at both fractions; one
step at batch 1 takes 8.3, 4.4 and 5.4\,ms.}
\label{tab:resources}
\end{table}

\begin{figure}[!htbp]
\centering
\figureorpending{fig_time.pdf}{width=\linewidth}
\caption{Where each run's GPU time went: active training (the 4-hour budget), waiting for data, and other overhead
(start-up, the copy of the data to the node's disk, validation, checkpoint writes, and steps replayed after a
preemption or crash). SU: service units estimated from the run logs at 40 per hour.}
\label{fig:time}
\end{figure}

\textbf{Compute for the whole project} (Table~\ref{tab:compute}). The project used 146.4 GPU-hours, 128.0 of them on
the paid partition (5,111 service units). The twelve core runs held 81.1 paid GPU-hours plus 2.2 on the free
partition, against the abstract's estimate of about 100 and the mid-project projection of 131--142; the copy of the
data to the node's disk saved the difference.

\begin{table}[!htbp]
\centering\footnotesize
\begin{tabularx}{\linewidth}{@{}>{\raggedright\arraybackslash}Xrrr@{}}
\toprule
Purpose & Jobs & GPU-hours & SU \\
\midrule
Core study: 10\% runs & 9 & 47.7 & 1,907 \\
Core study: 100\% runs & 12 & 33.4 & 1,337 \\
Test evaluations (core runs at 1 and 4 hours, persistence, released checkpoint, dedicated 1-hour runs) & 20 & 7.6 & 303 \\
Uncertainty: ensemble training (ten 1-hour runs) & 10 & 16.1 & 644 \\
Uncertainty and noise evaluation & 16 & 10.6 & 425 \\
Dedicated 1-hour MSE runs (point-forecast baseline for the ensemble) & 8 & 10.1 & 403 \\
Pipeline check & 1 & 2.3 & 84 \\
Numerical solver and flow snapshots (CPU work holding one GPU) & 4 & 0.2 & 7 \\
Free partitions (learning-rate search, profiling, first attempts, solver, data, setup) & 148 & 18.3 & 0 \\
\midrule
All & 228 & 146.4 & 5,111 \\
\bottomrule
\end{tabularx}
\caption{Compute billed for the whole project, from Slurm accounting. The paid partition bills 40 SU per hour for one
A100 and 8 cores; free partitions cost nothing but can be preempted.}
\label{tab:compute}
\end{table}

\section{Numerical-Solver Baseline}
\label{sec:solver}

A learned model is useful only if it is cheaper than solving the equations at the same accuracy~\cite{mcgreivy2024weak}.
The data were generated with Dedalus~\cite{burns2020dedalus} at $512\times128$, so we restart the same equations
(Dedalus, RK222 time stepping, 3/2 dealiasing) from the models' last input frame and run 30 frames ahead, on six
grids from $512\times128$ down to $64\times16$. The design is bounded: three regimes ($\mathrm{Ra}=10^6$, $10^8$ and
$10^{10}$ at $\mathrm{Pr}=1$), five test trajectories each and starts 80 and 140, 10 rollouts per regime, each on one
CPU core. The matching rule was fixed before any solver result: the coarsest grid whose mean VRMSE is at or below our
100\%, four-hour CNextU-Net's at steps 1, 10 and 30. Cost is compared in service-unit equivalents (one per core-hour,
32 per GPU-hour), and the network's training cost (185 SU per run) is added to find the number of forecasts after
which training pays off.

\begin{table}[!htbp]
\centering\scriptsize
\setlength{\tabcolsep}{3.5pt}
\begin{tabular}{@{}lcccccc@{}}
\toprule
& CNextU-Net (100\%, 4 h) & Matching & Solver & Seconds per rollout: & SU-eq.\ per rollout: & Break-even \\
Regime & steps 1 / 10 / 30 & grid & steps 1 / 10 / 30 & solver / model & solver / model & rollouts \\
\midrule
$\mathrm{Ra}=10^6$ & 0.011 / 0.112 / 0.355 & $128\times32$ & 0.0008 / 0.0010 / 0.0011 & 15.7 / 0.25 & 0.0044 / 0.0022 & $\sim$86,000 \\
$\mathrm{Ra}=10^8$ & 0.029 / 0.172 / 0.378 & $256\times64$ & 0.008 / 0.136 / 0.251 & 115 / 0.25 & 0.032 / 0.0022 & $\sim$6,200 \\
$\mathrm{Ra}=10^{10}$ & 0.099 / 0.329 / 0.618 & none & -- & 1,278$^{a}$ / 0.25 & -- & -- \\
\bottomrule
\end{tabular}
\caption{The coarsest solver grid as accurate as CNextU-Net at steps 1, 10 and 30, on the same 10 rollouts per
regime. Solver on one CPU core, model on one A100 at batch 1. Break-even: forecasts after which the model's training
cost plus its per-forecast cost falls below the solver's. $^{a}$The finest grid, $512\times128$.}
\label{tab:solver}
\end{table}

\begin{figure}[!htbp]
\centering
\figureorpending{fig_solver.pdf}{width=\linewidth}
\caption{The numerical solver against CNextU-Net (100\% data, four hours, seed mean) on the same 10 rollouts per
regime. Left three panels: VRMSE at each rollout step for the six solver grids (dark: fine; light: coarse),
CNextU-Net, the released CNextU-Net and persistence. Right: each grid's worst ratio to CNextU-Net's error over steps
1, 10 and 30 against its seconds per rollout; below 1 is as accurate as CNextU-Net. Open marker: diverged.}
\label{fig:solver}
\end{figure}

\textbf{Results} (Table~\ref{tab:solver}, Fig.~\ref{fig:solver}). At $\mathrm{Ra}=10^6$ the solver wins outright: from
$128\times32$ up it is more than 100 times more accurate than CNextU-Net at steps 10 and 30, at about twice the
model's cost per forecast in SU equivalents, so the model's error there is model error, not chaos. At $10^8$ the
$256\times64$ grid matches, at 115\,s per rollout, 465 times the model's wall time and 15 times its SU equivalent;
training pays off after about 6,200 forecasts. On wall time the model is 60 to 5,000 times faster, but a GPU-hour is
billed 32 times a core-hour, so per forecast the solver costs only 2--15 times as much before training is counted.
More data hardly moves the match: against our 10\% CNextU-Net the $\mathrm{Ra}=10^8$ match is $192\times48$ (break-even
after about 12,900 forecasts). Against the released CNextU-Net, which on these rollouts is more accurate one step ahead
at $10^8$ and $10^{10}$ but less accurate after 10 and 30 steps in every regime, the matches are $128\times32$,
$192\times48$ and none, as for our 10\% model.

\textbf{Limitation at $\mathrm{Ra}=10^{10}$.} No grid matches there (step 10: 0.37--0.64 against 0.33), and
$64\times16$ diverges in 3 of 10 rollouts. This result is affected by our time step. We used the Dedalus example's CFL
safety factor of 0.5; the data's generator used 0.0039, so whenever the CFL limit binds our steps are up to 128 times
larger. At $10^{10}$ our $512\times128$ restart already has VRMSE 0.167 after one step and piles up energy at the grid
scale (9 times the true energy at $k_x=120$--255), the signature of too large a step. So we report that no grid matches
CNextU-Net \emph{with our solver's time step}; a run with the data's step would cost up to about 45 hours per
$512\times128$ rollout. The matches at $10^6$ and $10^8$ stand: a smaller step would only make the solver more accurate
and more expensive. The cost of generating the training data with this solver is not counted.

\section{Robustness to Input Noise}
\label{sec:noise}

Measured or assimilated initial states are noisy, so we re-evaluated the four-hour checkpoints with Gaussian noise of
standard deviation $\sigma\in\{0.01, 0.05, 0.10, 0.20\}$ (in normalised units) added once to the four input frames,
with the same random draws for every model (evaluation only; Table~\ref{tab:noise}, Fig.~\ref{fig:noise}).

\begin{table}[!htbp]
\centering\small
\begin{tabular}{@{}llccccc@{}}
\toprule
& & \multicolumn{4}{c}{VRMSE relative to clean inputs} & VRMSE, \\
\cmidrule(lr){3-6}
Model (4 h) & Data & $\sigma=0.01$ & 0.05 & 0.10 & 0.20 & clean $\to$ $\sigma=0.20$ \\
\midrule
CNextU-Net & 10\% & 1.00 & 1.04 & 1.14 & \textbf{1.46} & 0.0505 $\to$ \textbf{0.0737} \\
           & 100\% & 1.00 & 1.04 & 1.15 & \textbf{1.47} & 0.0489 $\to$ \textbf{0.0719} \\
TFNO & 10\% & 1.01 & 1.09 & 1.25 & 1.57 & 0.1106 $\to$ 0.1740 \\
     & 100\% & 1.01 & 1.10 & 1.27 & 1.62 & 0.1075 $\to$ 0.1746 \\
AViT & 10\% & 1.02 & 1.33 & 1.89 & 2.94 & 0.0761 $\to$ 0.2238 \\
     & 100\% & 1.07 & 1.88 & 2.87 & 4.27 & 0.0636 $\to$ 0.2716 \\
\midrule
Persistence & -- & 1.00 & 1.07 & 1.20 & 1.56 & 0.2295 $\to$ 0.3585 \\
\bottomrule
\end{tabular}
\caption{One-step test VRMSE in developed convection with noisy inputs, seed mean. Bold: most robust.}
\label{tab:noise}
\end{table}

\begin{figure}[!htbp]
\centering
\figureorpending{fig_noise.pdf}{scale=1}
\caption{Test VRMSE against input noise: (a) one step in developed convection; (b) rollout step 10 from frames 80
and 140. Filled, solid: 100\% data; open, dashed: 10\%; dotted: persistence.}
\label{fig:noise}
\end{figure}

CNextU-Net is the most robust and stays the most accurate at every noise level; TFNO degrades about like
persistence. AViT is the most fragile, and more data makes it more so: the 100\% model is better on clean inputs (0.064
against 0.076) but worse from $\sigma=0.05$ (0.120 against 0.101), and from $\sigma=0.05$ on it falls behind TFNO. A
reading consistent with this, not tested: its patch embedding maps noise within a patch to tokens it never saw in
training, and more training tightens its fit to clean inputs. Over a rollout the input noise fades: at step 10 the
error rises by at most 13\% at $\sigma=0.2$ for CNextU-Net, TFNO and persistence (AViT 39--53\%), because the
rollout's own error dominates. Over all windows the ratios are much larger (up to 64 times for AViT and 39 for
persistence at $\sigma=0.2$): in the nearly still onset windows any noise is large relative to the true field.

\section{Uncertainty Quantification}
\label{sec:uq}

A forecast is more useful with an honest error bar. We asked how far each family's forecast can be trusted, and how
its uncertainty changes with data and along a rollout, in two parts.

\textbf{Methods.} (A) A \emph{deep ensemble}~\cite{lakshminarayanan2017ensembles} of five CNextU-Nets whose last layer
outputs a mean and a log-variance for every field and grid point, trained with the $\beta$-NLL loss
($\beta=1$)~\cite{seitzer2022pitfalls} for one hour at 10\% and at 100\% data (ten runs, otherwise the core recipe).
The forecast is the uniform mixture of the five Gaussians: its variance splits into the mean predicted variance
(\emph{aleatoric}) and the spread of the members' means (\emph{epistemic}). Two MSE-trained CNextU-Nets with the same
one-hour schedule serve as the point-forecast baseline. (B) For all three families, the two seeds of each four-hour
run form a two-member ensemble, and \emph{split conformal prediction}~\cite{angelopoulos2021gentle} turns any point
forecast into intervals with guaranteed average coverage: the score $|y-\hat y|/\sigma(x)$ is calibrated on the
validation split, per field and rollout step, and its quantile sets the interval width on test. The scale $\sigma(x)$
is constant (a plain point forecast), the members' spread, or the ensemble's predictive standard deviation. We report
coverage, interval width, the continuous ranked probability score (CRPS)~\cite{gneiting2007scoring}, the spread-skill
ratio (ensemble spread over error, corrected for ensemble size, so 1 is consistent), rank histograms and PIT
histograms~\cite{hamill2001rank}, all in normalised units at 128 random grid points per window. Three shifts test the
intervals: longer rollouts (the step-1 calibration kept), noisy inputs (Section~\ref{sec:noise}) and the highest
Rayleigh number held out of the calibration.

\begin{table}[!htbp]
\centering\small
\begin{tabular}{@{}lccc@{}}
\toprule
$\beta$-NLL ensemble (CNextU-Net, 1 h, 5 members) & One step, developed & Rollout step 10 & Rollout step 30 \\
\midrule
Ensemble-mean VRMSE, 10\% / 100\% data & 0.0792 / 0.0775 & 0.271 / 0.263 & 0.505 / 0.493 \\
CRPS (MSE model's absolute error), 100\% & 0.034 (0.047) & 0.133 (0.182) & 0.262 (0.388) \\
Coverage of the 50\% / 90\% intervals, 100\% & 0.71 / 0.95 & 0.39 / 0.73 & 0.42 / 0.76 \\
Spread-skill ratio, 100\% & 1.13 & 0.69 & 0.90 \\
Aleatoric share of the predicted variance & 91\% & 11\% & 2\% \\
\bottomrule
\end{tabular}
\caption{The deep ensemble with a variance output on the test split (rollouts from frames 80 and 140). The MSE pair's
ensemble-mean VRMSE is 0.0773 at one step at both fractions.}
\label{tab:uqa}
\end{table}

\textbf{The ensemble} (Table~\ref{tab:uqa}, Fig.~\ref{fig:uq}a). The variance output costs no accuracy: the
ensemble mean is as accurate as two MSE runs averaged, and its forecast distribution is 26--32\% better in CRPS than the
absolute error of an MSE point forecast at every horizon, and on every field (one step in developed convection 24--29\%,
least for pressure; during onset 16--21\%; at step 30 28--40\%, most for buoyancy). One step ahead its intervals are slightly too wide (96\% of
points inside the 90\% interval, spread-skill 1.13). Along a rollout they become too narrow (73\% by step 10): the
predicted variance shrinks on the model's own smoother inputs, while the members drift apart, so the epistemic part
grows thirtyfold and dominates from step 4. The ensemble predicts each step's uncertainty but does not propagate it.
After one hour the 10\% and 100\% ensembles agree to 1--3\% on every measure, so more data does not shrink their
uncertainty at this budget. Along the height, the predicted spread follows the error peak a few rows above the bottom
wall (16\% short), but on the wall rows themselves it is twice the error: the model does not learn that the wall values
are fixed. Across regimes, one step ahead, the predicted variance is 85--95\% aleatoric at every Rayleigh number. Both
parts grow with it from $10^6$ to $10^{10}$ (aleatoric 17-fold, epistemic 8-fold), in step with the squared error
(21-fold), so the ensemble is too wide at low Rayleigh number (spread-skill 1.30 in developed convection at $10^6$,
1.45 during onset) and nearly consistent at $10^{10}$ (1.10).

\textbf{Seed pairs and conformal intervals} (Table~\ref{tab:conformal}). The two seeds of a four-hour run agree far
more closely than either agrees with the truth (spread-skill 0.41--0.70 one step ahead, rising to 0.84--0.89 by step
30), although their rank histograms are nearly flat; a flat rank histogram does not imply
calibration~\cite{hamill2001rank}. They are most under-dispersed at high Rayleigh number (spread-skill in developed
convection 0.69--1.16 at $10^6$, 0.34--0.67 at $10^{10}$). More data makes the seeds agree more: their one-step epistemic variance falls by
71\% for AViT, 35\% for CNextU-Net and 15\% for TFNO, the data effect the one-hour ensembles do not show. Split-conformal
90\% intervals hold their coverage on test for every model and scale (0.899--0.904), so validation and test behave as
exchangeable; the coverage is an average over phases (onset 0.94--0.99, developed 0.87--0.88). Their widths rank the
models as their errors do, and the ensemble's predictive scale gives intervals 30\% narrower than its MSE counterpart
one step ahead.

\begin{table}[!htbp]
\centering\small
\begin{tabular}{@{}lcc@{}}
\toprule
Width of the 90\% conformal interval (normalised) & One step, 10\% / 100\% & Rollout step 30, 10\% / 100\% \\
\midrule
CNextU-Net, 4 h & 0.116 / 0.112 & 1.74 / 1.71 \\
TFNO, 4 h & 0.293 / 0.284 & 1.73 / 1.67 \\
AViT, 4 h & 0.185 / 0.153 & 2.39 / 2.49 \\
Persistence & 0.654 & 2.49 \\
\midrule
CNextU-Net $\beta$-NLL ensemble, 1 h (predictive scale) & 0.140 / 0.136 & 1.64 / 1.61 \\
CNextU-Net MSE, 1 h (constant scale) & 0.199 / 0.197 & 1.95 / 1.99 \\
\bottomrule
\end{tabular}
\caption{Width of split-conformal 90\% intervals on test (field mean), calibrated on validation per step; every row
covers 0.899--0.904 on test. Constant scale unless stated.}
\label{tab:conformal}
\end{table}

\textbf{Shifts} (Fig.~\ref{fig:uq}b). (a) \emph{Longer rollouts}: an interval of constant width calibrated at step 1
collapses along the rollout (coverage at step 30: CNextU-Net 0.31, TFNO 0.43, AViT 0.24), while an interval scaled by
the members' spread stays valid at every step (seed pairs 0.87--0.91, the five-member ensemble 0.92--0.94). The
members' disagreement tracks how the rollout error grows; the predicted variance does not. (b) \emph{Noisy inputs}:
at $\sigma=0.2$ coverage falls to 0.76--0.78 for CNextU-Net and 0.70--0.73 for TFNO, and to 0.42--0.56 for AViT, which
the spread scale lifts to 0.81--0.86. (c) \emph{Rayleigh number}: calibrated on $\mathrm{Ra}\le10^9$, coverage at
$10^{10}$ is 0.66 for CNextU-Net, 0.70 for AViT and 0.76 for TFNO, 0.85 with the ensemble's predictive scale;
recalibrating on just two $\mathrm{Ra}=10^{10}$ validation trajectories restores 0.84--0.90.

\begin{figure}[!htbp]
\centering
\figureorpending{fig_uq.pdf}{width=\linewidth}
\caption{Calibration along rollouts from frames 80 and 140 (test). (a) Coverage of the five-member ensemble's
Gaussian-mixture 50\% and 90\% intervals. (b) Coverage of split-conformal 90\% intervals calibrated at step 1 and kept,
four-hour seed pairs at 100\% data: constant width (dashed) against spread-scaled (solid). (c) Spread-skill ratio of
the seed pairs (solid: 100\%, dashed: 10\%) and the ensemble (thick).}
\label{fig:uq}
\end{figure}

\textbf{Reliability and coverage window by window} (Fig.~\ref{fig:uqrel}, Table~\ref{tab:uqcond}). An average coverage
of 90\% can hide where intervals fail, so we also binned the error by the predicted standard deviation and computed
each window's own coverage. (a) \emph{Reliability.} The ensemble's predicted standard deviation carries real
information: one step ahead the RMSE rises 46- to 54-fold from its lowest to its highest decile. In the top decile the
RMSE is close to the predicted value (0.93--0.94 of it), in the lowest about half of it, so the intervals are too wide
where the flow is easy. By step 10 the RMSE is 1.3--1.7 times the predicted value in every decile, the narrowness seen
above. The seed pairs' spread also ranks the errors (one step ahead the RMSE rises 8- to 12-fold over the deciles for
CNextU-Net and AViT, 5- to 6-fold for TFNO), but with two members it is often near zero by chance, so its lowest
deciles lie far above the diagonal. (b) \emph{Coverage by window and Rayleigh number.} Constant-width conformal
intervals reach their 90\% only on average: they over-cover low Rayleigh numbers (0.97--0.99 at $10^6$) and
under-cover high ones (0.76--0.85 at $10^{10}$), and about a fifth of the windows are covered below 80\%. Scaling by
the seed spread or by the ensemble's predictive standard deviation, without new training, makes coverage much more even
across windows (1--4\% of windows below 80\%) and across Rayleigh numbers (0.92--0.94 at $10^6$, 0.84--0.87 at
$10^{10}$). Recalibrating on two trajectories at the highest Rayleigh number (shift c) closes most of the remaining
gap. Along rollouts the pattern holds:
at step 10, 11--15\% of the windows fall below 80\% coverage with constant width, at most 3\% with a spread scale.

\begin{table}[!htbp]
\centering\footnotesize
\setlength{\tabcolsep}{4pt}
\begin{tabularx}{\linewidth}{@{}>{\raggedright\arraybackslash}Xccccc@{}}
\toprule
& \multicolumn{3}{c}{One step} & \multicolumn{2}{c}{Rollout step 10} \\
\cmidrule(lr){2-4}\cmidrule(l){5-6}
& Coverage at & \multicolumn{2}{c}{Windows' own coverage} & Coverage at & Windows \\
90\% interval (100\% data) & $\mathrm{Ra}=10^6$ / $10^{10}$ & 10th pct. & below 0.8 & $\mathrm{Ra}=10^6$ / $10^{10}$ & below 0.8 \\
\midrule
CNextU-Net, 4 h: constant width & 0.99 / 0.76 & 0.68 & 19\% & 0.96 / 0.84 & 15\% \\
\quad scaled by the seed spread & 0.93 / 0.84 & 0.82 & 4\% & 0.91 / 0.87 & 1\% \\
TFNO, 4 h: constant width & 0.99 / 0.81 & 0.69 & 20\% & 0.96 / 0.85 & 12\% \\
\quad scaled by the seed spread & 0.94 / 0.85 & 0.82 & 4\% & 0.93 / 0.87 & 1\% \\
AViT, 4 h: constant width & 0.99 / 0.78 & 0.67 & 20\% & 0.93 / 0.89 & 11\% \\
\quad scaled by the seed spread & 0.94 / 0.84 & 0.82 & 4\% & 0.93 / 0.87 & 0\% \\
Persistence: constant width & 0.97 / 0.85 & 0.70 & 21\% & 0.90 / 0.90 & 15\% \\
\midrule
$\beta$-NLL ensemble, 1 h: conformal, predictive scale & 0.92 / 0.87 & 0.83 & 1\% & 0.91 / 0.88 & 3\% \\
\quad Gaussian mixture, no calibration & 0.97 / 0.95 & 0.93 & 0\% & 0.74 / 0.69 & 84\% \\
\bottomrule
\end{tabularx}
\caption{Coverage of 90\% intervals on test beyond the average (field mean): at the lowest and highest Rayleigh number,
and the distribution of each window's own coverage (128 grid points $\times$ 4 fields per window; 8,575 one-step
windows, 350 rollouts from frames 80 and 140). Split-conformal rows are calibrated on validation (one step pooled,
rollouts per step); every one averages 0.90 on test. The 10\% runs give the same picture.}
\label{tab:uqcond}
\end{table}

\begin{figure}[!htbp]
\centering
\figureorpending{fig_uq_reliability.pdf}{width=\linewidth}
\caption{Reliability and coverage by Rayleigh number, 100\% data. (a, b) RMSE of the ensemble mean against its
predicted standard deviation, in deciles of the predicted value (four fields pooled, normalised units): four-hour seed
pairs (spread with the finite-ensemble factor) and the one-hour $\beta$-NLL ensemble (thick), one step ahead and at
rollout step 10 from frames 80 and 140; on the diagonal the predicted value is right. (c) Coverage of split-conformal
90\% intervals one step ahead per Rayleigh number: constant width (dashed) against scaled by the seed spread (solid)
or the ensemble's predictive standard deviation (thick); persistence dotted.}
\label{fig:uqrel}
\end{figure}

\section{Findings Under Limited Data and Compute}
\label{sec:findings}

No family is best on every count (Table~\ref{tab:choice}). CNextU-Net is the most accurate and stable forecaster of
developed convection and the most robust to noisy inputs, TFNO the fastest for single forecasts and the closest in heat
transport, and AViT the cheapest on the GPU and the one that gains most from more data.

\begin{table}[!htbp]
\centering\footnotesize
\begin{tabularx}{\linewidth}{@{}>{\raggedright\arraybackslash}p{0.25\linewidth}>{\raggedright\arraybackslash}p{0.2\linewidth}>{\raggedright\arraybackslash}X@{}}
\toprule
If the priority is & Best of the three & Trade-off and evidence \\
\midrule
Accurate forecasts of developed convection that stay stable & CNextU-Net, at both fractions & One step 0.0505
(10\%) / 0.0489 (100\%); 88--89\% of rollouts stay below VRMSE 1; slowest inference (248\,ms per 30 steps); 61\,GB
training memory. 1.28 times the released model's one-step error, but more stable over rollouts \\ \addlinespace[3pt]
Fastest single forecast (batch 1) & TFNO (133--135\,ms per 30 steps) & Largest one-step error (0.107--0.111); too
smooth from $k_x\approx25$; gains least from longer training; unstable rollouts after one hour at 100\% \\
\addlinespace[3pt]
Least training memory, cheapest batched forecasts & AViT (8.7\,GB; 0.89\,ms per sample at batch 64) & Rollouts
degrade fastest (behind persistence after 21--25 steps); patch-shaped artifacts; most fragile to input noise \\
\addlinespace[3pt]
More training data at the same budget & AViT & $-16.5$\% one-step error from 10\% to 100\% (CNextU-Net $-3.1$\%, TFNO
$-2.8$\%); no model's rollouts improve; one to four hours matters more for CNextU-Net ($-35$\%) \\ \addlinespace[3pt]
Only one hour of training & 10\%: AViT (one step), TFNO (rollouts). 100\%: CNextU-Net or AViT & 100\%: one step
0.456 for both, developed 0.075 / 0.080; 50--52\% of rollouts below VRMSE 1, against 23\% for TFNO \\
\addlinespace[3pt]
Convection onset & None: persistence & Persistence 0.116 against 0.245--0.452 at four hours; the deficit is a velocity
effect, on buoyancy all three are more accurate; every model is better from about 20 frames before each
trajectory's onset frame \\ \addlinespace[3pt]
Heat transport over a rollout & TFNO & Median error 20--21\%; AViT carries too little heat ($-17$\%); CNextU-Net
scattered (49\%); the released CNextU-Net twice the true value \\ \addlinespace[3pt]
Noisy inputs & CNextU-Net & Error $\times1.47$ at $\sigma=0.2$, against $\times1.62$ (TFNO) and $\times4.27$ (AViT,
100\%) \\ \addlinespace[3pt]
Calibrated uncertainty & CNextU-Net ensemble with conformal intervals & CRPS 26--32\% below a point forecast's error;
intervals too narrow along rollouts unless scaled by the members' spread (87--94\% at step 30), which also evens
coverage across windows and Rayleigh numbers; five training runs \\
\addlinespace[3pt]
Many forecasts at $\mathrm{Ra}\le10^8$ & CNextU-Net, above a break-even & A coarse solver grid is as accurate
(more than 100 times more accurate at $10^6$); training pays off after $\sim$86,000 ($10^6$) and $\sim$6,200 ($10^8$)
forecasts \\
\bottomrule
\end{tabularx}
\caption{Which model to use under which constraint, after four hours of training unless the row says otherwise. These
hold for these configurations and budgets.}
\label{tab:choice}
\end{table}

\textbf{Negative results.} (a) Ten times more data does not improve any model's 30-step rollouts at these budgets,
and it does not close the gap to the longer-trained released CNextU-Net. (b) A low one-step error does not guarantee
a stable rollout or the right heat transport: the released CNextU-Net has the lowest one-step error in developed
convection and the least accurate heat transport. (c) No model beats persistence on average over onset, although every model is better from about 20
frames before each trajectory's onset frame. (d) At
$\mathrm{Ra}=10^6$ a coarse numerical solver is far more accurate than every model. (e) The ensemble's predicted
variance does not track rollout error, and more data makes AViT more fragile to noise.

\textbf{Limitations.} (a) Two seeds per configuration, both on one fixed 10\% subset, so the effect of which
trajectories were drawn is not measured. (b) Training is one-step only; multi-step training might change the rollout
results~\cite{brandstetter2022mp}. (c) Training in fp32 favours TFNO, since bf16 would speed up CNextU-Net 1.7 times
and AViT 2.2 times but is not available for TFNO. (d) The learning rate is the best of a short search and lies at the
top of the tested grid. (e) The solver comparison at $\mathrm{Ra}=10^{10}$ is limited by our time step
(Section~\ref{sec:solver}). (f) The one-hour results are checkpoints of four-hour runs, not dedicated one-hour
schedules, except for the uncertainty ensemble.

\section{Conclusion}
\label{sec:conclusion}

Under limited data and compute, the choice of architecture matters more than ten times more training data. With one
recipe and four hours on one A100, CNextU-Net is the most accurate and most stable forecaster of developed
Rayleigh--B\'enard convection at both data fractions, and the most robust to noisy inputs; more data barely helps it,
while more training time does. AViT gains most from data but its rollouts degrade fastest and it is fragile to noise;
TFNO is the fastest to run and carries the right heat transport, but is the least accurate pointwise. None beats
copying the last frame during onset, and at low Rayleigh numbers a coarse numerical solver matches the networks for
the cost of tens of thousands of forecasts. A small ensemble with conformal intervals scaled by its spread gives
error bars that stay valid along a rollout.

\clearpage

\phantomsection\addcontentsline{toc}{section}{References}
\begin{thebibliography}{99}
\setlength{\itemsep}{8pt}

\bibitem{ohana2024well}
Ohana, R., et al.
``The Well: a Large-Scale Collection of Diverse Physics Simulations for Machine Learning.''
NeurIPS, Datasets and Benchmarks Track, 2024.\\
\url{https://arxiv.org/abs/2412.00568}

\bibitem{liu2022convnext}
Liu, Z., Mao, H., Wu, C.-Y., Feichtenhofer, C., Darrell, T., and Xie, S.
``A ConvNet for the 2020s.''
CVPR, 2022.\\
\url{https://arxiv.org/abs/2201.03545}

\bibitem{kossaifi2024tfno}
Kossaifi, J., Kovachki, N., Azizzadenesheli, K., and Anandkumar, A.
``Multi-Grid Tensorized Fourier Neural Operator for High-Resolution PDEs.''
TMLR, 2024.\\
\url{https://arxiv.org/abs/2310.00120}

\bibitem{mccabe2024mpp}
McCabe, M., et al.
``Multiple Physics Pretraining for Spatiotemporal Surrogate Models.''
NeurIPS, 2024.\\
\url{https://arxiv.org/abs/2310.02994}

\bibitem{li2021fno}
Li, Z., Kovachki, N., Azizzadenesheli, K., Liu, B., Bhattacharya, K., Stuart, A., and Anandkumar, A.
``Fourier Neural Operator for Parametric Partial Differential Equations.''
ICLR, 2021.\\
\url{https://arxiv.org/abs/2010.08895}

\bibitem{rahaman2019spectral}
Rahaman, N., Baratin, A., Arpit, D., Draxler, F., Lin, M., Hamprecht, F. A., Bengio, Y., and Courville, A.
``On the Spectral Bias of Neural Networks.''
ICML, 2019.\\
\url{https://arxiv.org/abs/1806.08734}

\bibitem{lippe2023pderefiner}
Lippe, P., Veeling, B. S., Perdikaris, P., Turner, R. E., and Brandstetter, J.
``PDE-Refiner: Achieving Accurate Long Rollouts with Neural PDE Solvers.''
NeurIPS, 2023.\\
\url{https://arxiv.org/abs/2308.05732}

\bibitem{brandstetter2022mp}
Brandstetter, J., Worrall, D., and Welling, M.
``Message Passing Neural PDE Solvers.''
ICLR, 2022.\\
\url{https://arxiv.org/abs/2202.03376}

\bibitem{mccabe2023stability}
McCabe, M., Harrington, P., Subramanian, S., and Brown, J.
``Towards Stability of Autoregressive Neural Operators.''
TMLR, 2023.\\
\url{https://arxiv.org/abs/2306.10619}

\bibitem{burns2020dedalus}
Burns, K. J., Vasil, G. M., Oishi, J. S., Lecoanet, D., and Brown, B. P.
``Dedalus: A flexible framework for numerical simulations with spectral methods.''
Physical Review Research 2, 023068, 2020.\\
\url{https://arxiv.org/abs/1905.10388}

\bibitem{mcgreivy2024weak}
McGreivy, N., and Hakim, A.
``Weak baselines and reporting biases lead to overoptimism in machine learning for fluid-related partial
differential equations.''
Nature Machine Intelligence 6, 2024.\\
\url{https://arxiv.org/abs/2407.07218}

\bibitem{lakshminarayanan2017ensembles}
Lakshminarayanan, B., Pritzel, A., and Blundell, C.
``Simple and Scalable Predictive Uncertainty Estimation using Deep Ensembles.''
NeurIPS, 2017.\\
\url{https://arxiv.org/abs/1612.01474}

\bibitem{seitzer2022pitfalls}
Seitzer, M., Tavakoli, A., Antic, D., and Martius, G.
``On the Pitfalls of Heteroscedastic Uncertainty Estimation with Probabilistic Neural Networks.''
ICLR, 2022.\\
\url{https://arxiv.org/abs/2203.09168}

\bibitem{angelopoulos2021gentle}
Angelopoulos, A. N., and Bates, S.
``A Gentle Introduction to Conformal Prediction and Distribution-Free Uncertainty Quantification.''
arXiv:2107.07511, 2021.\\
\url{https://arxiv.org/abs/2107.07511}

\bibitem{gneiting2007scoring}
Gneiting, T., and Raftery, A. E.
``Strictly Proper Scoring Rules, Prediction, and Estimation.''
Journal of the American Statistical Association 102(477), 359--378, 2007.\\
\url{https://doi.org/10.1198/016214506000001437}

\bibitem{hamill2001rank}
Hamill, T. M.
``Interpretation of Rank Histograms for Verifying Ensemble Forecasts.''
Monthly Weather Review 129(3), 550--560, 2001.

\bibitem{wellfix}
The Well source code, commits
\href{https://github.com/PolymathicAI/the_well/commit/2d8d7eaf9eb46b96f215abf39c94b6be4bf9ea5c}{\texttt{2d8d7eaf9eb46b96f215abf39c94b6be4bf9ea5c}}
(2024-06-03) and
\href{https://github.com/PolymathicAI/the_well/commit/e4301602e57f2dddeaa87856dda49bbba47c663b}{\texttt{e4301602e57f2dddeaa87856dda49bbba47c663b}}
(2025-11-17).\\
\url{https://github.com/PolymathicAI/the_well/commit/2d8d7eaf9eb46b96f215abf39c94b6be4bf9ea5c}\\
\url{https://github.com/PolymathicAI/the_well/commit/e4301602e57f2dddeaa87856dda49bbba47c663b}

\bibitem{wellrb}
The Well, \texttt{rayleigh\_benard} dataset card, commit
\href{https://github.com/PolymathicAI/the_well/commit/e63dec8ded515be864755eaf595d39fea409e335}{\texttt{e63dec8}}.\\
\url{https://github.com/PolymathicAI/the_well/blob/e63dec8/datasets/rayleigh_benard/README.md}

\end{thebibliography}

\clearpage
\appendix
\addtocontents{toc}{\protect\appendixcontents}
% Appendix headings read "Appendix A: ..."; \ref still gives the bare letter, so "Appendix~\ref{...}" is unchanged.
\makeatletter
\renewcommand{\@seccntformat}[1]{\ifcsname prefix@#1\endcsname\csname prefix@#1\endcsname\fi
  \csname the#1\endcsname\ifcsname prefix@#1\endcsname:\fi\quad}
\newcommand{\prefix@section}{Appendix~}
\makeatother

\section{Physical System Details}
\label{app:rb_regimes}

The data solve the two-dimensional Boussinesq equations with no-slip walls~\cite{wellrb}. A regime is set by
\begin{equation}
    \mathrm{Ra} = \frac{g\alpha\,\Delta T\,H^3}{\nu\kappa}, \qquad \mathrm{Pr} = \frac{\nu}{\kappa},
\end{equation}
with gravity $g$, thermal expansion coefficient $\alpha$, wall temperature difference $\Delta T$, layer height $H$,
viscosity $\nu$ and thermal diffusivity $\kappa$; $\mathrm{Pr}$ is about 0.7 for air and 7 for water. In the data's
units $\kappa=(\mathrm{Ra}\,\mathrm{Pr})^{-1/2}$ and $\nu=(\mathrm{Ra}/\mathrm{Pr})^{-1/2}$. The dataset crosses
$\mathrm{Ra}\in\{10^6,10^7,10^8,10^9,10^{10}\}$ with $\mathrm{Pr}\in\{0.1,0.2,0.5,1,2,5,10\}$, and the 50
trajectories of each regime cross five initial-profile amplitudes $\delta b_0\in\{0.2,0.4,0.6,0.8,1.0\}$ with ten
noise seeds. The files store an evenly spaced placeholder for $y$; the values lie on the 128 Chebyshev--Gauss nodes
of the solver, the only grid on which the data meet their boundary conditions ($b=1$ and 0 at the walls within
$6\times10^{-10}$). The spectrum and the Nusselt number use these nodes; the models and VRMSE treat the rows as an
image.

\section{Released TFNO and Forward Pass}
\label{app:tfno}
In The Well's code, commit
\href{https://github.com/PolymathicAI/the_well/commit/2d8d7eaf9eb46b96f215abf39c94b6be4bf9ea5c}{\texttt{2d8d7eaf}}
(2024-06-03) called TFNO's Fourier blocks without keeping their output, so the model computed only its lifting and
projection layers; commit
\href{https://github.com/PolymathicAI/the_well/commit/e4301602e57f2dddeaa87856dda49bbba47c663b}{\texttt{e4301602}}
(2025-11-17) fixed it~\cite{wellfix}. The paper's TFNO results and the released checkpoint (2025-03-28) predate the
fix: with the forward pass it was trained with, the released TFNO reproduces the published value, and with the fixed
one it scores $\approx$10.5. Our runs use the fixed code (commit
\href{https://github.com/PolymathicAI/the_well/commit/e63dec8ded515be864755eaf595d39fea409e335}{\texttt{e63dec8}}).

\begin{table}[!htbp]
\centering\footnotesize
\begin{tabular}{@{}lccc@{}}
\toprule
& Paper-protocol one-step & Physical-field one-step & Rollout from frame 0, \\
Model & VRMSE (published) & all / onset / developed & steps 6--12 / 13--30 \\
\midrule
Released TFNO (paper-era forward) & 0.6634 (0.6566) & 0.300 / 0.594 / 0.193 & 142 / 169 (published $>10$) \\
Our TFNO (10\% data, 4 h) & 0.488 & 0.209 / 0.480 / 0.111 & 51 / 77 \\
Our TFNO (100\% data, 4 h) & 0.459 & 0.199 / 0.452 / 0.107 & 44 / 63 \\
Persistence & 0.202 & 0.199 / 0.116 / 0.229 & 0.27 / 0.50 \\
\bottomrule
\end{tabular}
\caption{The released TFNO, scored by our pipeline on the test split, against our TFNO (seed mean) and
persistence.}
\label{tab:released_tfno}
\end{table}

Without its Fourier blocks the model acts on each pixel separately and cannot move structures such as plumes, so it
learns little more than copying each pixel's last value (Fig.~\ref{fig:released}). It is slightly better than
persistence in developed convection (0.193 against 0.229) and much worse during onset. Our TFNO, with its Fourier
layers, is better everywhere at both data fractions.

\section{Metric Details}
\label{app:metrics}

\textbf{Physical-field VRMSE and the paper protocol.} VRMSE divides the error by how much the true field varies
across the grid. During onset the velocity barely varies, so tiny errors give large scores, held down only by
$\varepsilon=10^{-7}$ (Table~\ref{tab:units}, bottom). The paper protocol first divides each field by its training
standard deviation $\sigma_f$ (0.249, 0.149, 0.208 and 0.180 for $b$, $p$, $u$ and $w$); we score the fields as
stored. The rescaling cancels out except for $\varepsilon$, which becomes 16--45 times smaller relative to the
fields and so barely holds the onset scores down. In developed convection $\varepsilon$ is negligible and the two
agree. As a result onset, 27\% of the windows, gives 96\% of the released CNextU-Net's paper-protocol mean
(Table~\ref{tab:units}), and that mean can reverse rankings: at four hours AViT looks better than CNextU-Net (0.298
against 0.349 at 10\%; 0.245 against 0.333 at 100\%) but is 1.5 and 1.3 times less accurate in developed convection.
The physical-field VRMSE reduces this effect (87\%) without removing it, so we also report the two phases separately.

\begin{table}[!htbp]
\centering\small
\begin{tabular}{@{}lcc@{}}
\toprule
One-step windows & Physical-field VRMSE & Paper protocol \\
\midrule
Onset (27\% of the windows) & 0.733 & 2.400 \\
Developed convection & 0.0382 & 0.0382 \\
All windows & 0.224 & 0.669 (published 0.6699) \\
\midrule
Share of the all-window mean that comes from onset & 87\% & 96\% \\
\midrule
Velocity $u$ / $w$, frames 0--9 (onset) & 2.84 / 3.42 & 12.49 / 15.17 \\
Velocity $u$ / $w$, developed convection & 0.029 / 0.033 & 0.029 / 0.033 \\
\bottomrule
\end{tabular}
\caption{The released CNextU-Net's one-step test VRMSE scored both ways: four-field mean (top) and velocity alone
(bottom). The predictions are the same; only the units of the fields differ.}
\label{tab:units}
\end{table}

\textbf{Onset frame.} The kinetic energy is the domain mean of $(u^2+w^2)/2$ (uniform in $x$, Chebyshev--Gauss
quadrature in $y$), computed from the true frames of each trajectory. $t_{\mathrm{on}}$ is the first frame at which it
reaches 10\% of its mean over frames 150--199; every test trajectory crosses. The fraction was chosen on validation:
the released model's error is inflated only while the kinetic energy is below about 1--3\% of its developed mean, so a
fraction of 0.5 or more would count the growth of convection, where the flow already has structure, as onset. At 0.1
onset also ends later at low Rayleigh number, as expected physically; at 0.5 the onset frames hardly depend on it.

\textbf{Rollouts.} The paper's rollouts start at frame 0, so they lie entirely in onset; it reports the mean
normalised VRMSE over steps 6--12 and 13--30 (``$>10$'' for every model). Ours start at frames 20, 80 and 140; we
report both (Table~\ref{tab:rollout}).

\textbf{Kinetic-energy spectrum.} $E(k_x)=\tfrac12(|\hat u|^2+|\hat w|^2)$ from the Fourier transform of each row
along the periodic $x$ axis ($k_x=2\pi n/L_x$, $L_x=4$), averaged over rows with Chebyshev quadrature weights and
over each regime's rollouts at steps 1, 10 and 30. The log spectral error is the mean absolute $\log_{10}$ ratio of
predicted to true energy.

\textbf{Nusselt number.} $\mathrm{Nu}=1+\langle w\,b\rangle/\kappa = 1+\sqrt{\mathrm{Ra}\,\mathrm{Pr}}\,\langle
w\,b\rangle$ in the data's units, the volume average over the domain, computed for every predicted and true frame. The
factor $\sqrt{\mathrm{Ra}\,\mathrm{Pr}}$ (up to $3\times10^5$) amplifies small errors, and the instantaneous value
swings widely, so we average it over each 30-step rollout and report medians. Over 30 steps the true flow is not
statistically stationary (the heat flux through the bottom wall exceeds the volume average, which exceeds the top
wall's), so the values compare prediction with truth but are not meant for comparison with Nusselt--Rayleigh scaling
laws.

\textbf{Inference time.} Batch 1, fp32, one A100, input already on the GPU: after 10 warm-up predictions we time 50
one-step predictions and one 30-step rollout, each ending with a CUDA synchronisation.

\textbf{Uncertainty measures.} Coverage is the share of test values inside an interval. CRPS of the Gaussian mixture
is computed in closed form per point; for a point forecast it equals the absolute error. The spread-skill ratio is
$\sqrt{\overline{c\,\mathrm{var}_{\mathrm{ep}}+\mathrm{var}_{\mathrm{al}}}}\,/\,\mathrm{RMSE}$ of the ensemble mean,
with $c=(M+1)/(M-1)$ for $M$ members, so a consistent finite ensemble scores 1. Split conformal uses the
$\lceil (n+1)(1-\alpha)\rceil$-th smallest of $n$ validation scores; the spread scale is floored at the validation
set's 5th percentile. For reliability, every sampled grid point gets the predicted standard deviation
$\sqrt{c\,\mathrm{var}_{\mathrm{ep}}+\mathrm{var}_{\mathrm{al}}}$; the points are split into deciles of it, and each
decile's RMSE of the ensemble mean is compared with its root-mean-square predicted value. A window's own coverage is
the share of its 128 sampled points and four fields inside the interval.

\textbf{End of the onset error.} Per Rayleigh number, the first window start, counted from frame 0 or from each
trajectory's $t_{\mathrm{on}}$, at which the mean one-step VRMSE over the regime's trajectories falls below twice its
median in developed convection.

\section{Detailed Training Recipe}
\label{app:recipe}

Table~\ref{tab:recipe} compares our recipe with The Well's, taken from its paper and, where the two disagree, its
training configuration.

\textbf{Learning-rate search.} Each model was trained for 20 minutes at each rate (seed 0, 10\% data, nine runs),
with 100 warmup steps and 512 validation windows, and takes the rate with the lowest validation VRMSE
(Table~\ref{tab:lr}); the test split was not used. $10^{-3}$ won for all three.

\begin{table}[!htbp]
\centering\small
\begin{tabular}{@{}lccc@{}}
\toprule
Validation VRMSE after 20 min & $10^{-4}$ & $5\cdot10^{-4}$ & $10^{-3}$ \\
\midrule
CNextU-Net & 0.841 & 0.469 & \textbf{0.408} \\
TFNO & 0.603 & 0.358 & \textbf{0.302} \\
AViT & 0.646 & 0.332 & \textbf{0.257} \\
\bottomrule
\end{tabular}
\caption{Learning-rate search.}
\label{tab:lr}
\end{table}

\begin{table}[!htbp]
\centering\footnotesize
\begin{tabularx}{\linewidth}{@{}>{\raggedright\arraybackslash}p{2.0cm}>{\raggedright\arraybackslash}X>{\raggedright\arraybackslash}X@{}}
\toprule
& This study & The Well paper / code \\
\midrule
Activation & GELU (all three models) & GELU (same model code) \\
Optimizer & AdamW, $\beta = (0.9, 0.999)$, weight decay $10^{-4}$, no clipping & AdamW, weight decay $10^{-4}$ in the code (0.01 in the text), no clipping \\
Loss & one-step MSE, normalised units, equal field weights ($\beta$-NLL for the uncertainty ensemble) & one-step MSE \\
Batch size, precision & 64, fp32 & 64, fp32 \\
Learning rate & $10^{-3}$ for all, from our search over $\{10^{-4}, 5\cdot10^{-4}, 10^{-3}\}$ & per model from $\{10^{-4}, 5\cdot10^{-4}, 10^{-3}, 5\cdot10^{-3}, 10^{-2}\}$: $5\cdot10^{-4}$ (CNextU-Net), $10^{-4}$ (TFNO) \\
Schedule & 500 warmup steps, cosine to 0 over the 4-hour budget (1 hour for the ensemble) & 5 warmup epochs, cosine to $0.1\cdot$lr over 200 epochs, stopped at 12\,h \\
Budget, hardware & 1\,h and 4\,h of active GPU training on one A100 & 12\,h on one H100, best validation epoch \\
Training data & 10\% and 100\% of the trajectories & 100\% \\
\bottomrule
\end{tabularx}
\caption{Training recipe: ours against the paper's.}
\label{tab:recipe}
\end{table}

\textbf{Task and loss.} The model sees four consecutive frames and predicts the fifth. The fields are normalised
with the statistics of The Well's full training split, and the loss is the MSE over all grid points with equal
field weights. An epoch is 27,440 windows at 10\% data and 274,400 at 100\%.

\textbf{Checkpoints and seeds.} We evaluate the 1-hour and 4-hour checkpoints without selecting by validation score
(the paper kept its best validation epoch); 64 fixed validation windows are scored every 30 minutes for monitoring
only (Fig.~\ref{fig:curves}). A seed sets the initial weights and the data order; both seeds use the same subset. At
the 1-hour checkpoint the learning rate is still about 85\% of its peak.

\textbf{Software.} PyTorch 2.8.0, CUDA 12.8, Python 3.13; The Well's model code, with TFNO's layers from
neuraloperator 0.3.0; Dedalus 3.0.5 for the solver baseline.

\textbf{Reproducibility.} The code, configurations, Slurm scripts, run manifests and tests are in the project repository
(\url{https://github.com/skagankose/sciml-project}). Each run's configuration is hashed into its run identifier, and
production runs refuse to start from an uncommitted code tree. The pinned upstream versions are The Well
\texttt{e63dec8}, neuraloperator 0.3.0 and PyTorch 2.8.0. From the evaluated outputs, \texttt{rbbench aggregate},
\texttt{scripts/findings.py} and \texttt{scripts/final\_figures.py} rebuild every table and figure of this report.

\section{Resource Accounting}
\label{app:accounting}

\begin{table}[!htbp]
\centering\footnotesize
\begin{tabularx}{\linewidth}{@{}>{\raggedright\arraybackslash}X*{3}{>{\centering\arraybackslash}p{2.3cm}}@{}}
\toprule
Hours per run & CNextU-Net & TFNO & AViT \\
\midrule
\multicolumn{4}{@{}l}{\textit{10\% runs, data read from the shared file system}} \\
GPU training (the budget) & 4.0 & 4.0 & 4.0 \\
Waiting for data & 0.8--1.1 & 2.5--2.8 & 6.1--7.3 \\
Other (start-up, validation, checkpoints, replayed steps) & 0.4--0.9 & 0.3 & 0.9--1.0 \\
\textbf{Allocated time of training} & \textbf{5.2--6.0} & \textbf{6.9--7.1} & \textbf{11.1--12.2} \\
\midrule
\multicolumn{4}{@{}l}{\textit{100\% runs, data copied to the GPU node's disk}} \\
GPU training (the budget) & 4.0 & 4.0 & 4.0 \\
Waiting for data & 0.0 & 0.0--0.1 & 0.2--0.9 \\
Other (start-up, data copy, validation, checkpoints, replayed steps) & 0.4--0.6 & 1.5--1.7 & 0.8--1.8 \\
\textbf{Allocated time of training} & \textbf{4.4--4.6} & \textbf{5.5--5.8} & \textbf{5.7--6.0} \\
\bottomrule
\end{tabularx}
\caption{Time per run (range over the two seeds). Queue waits before a job starts depend on the cluster's load and are
not included (up to 12--15 hours on the paid partition).}
\label{tab:time}
\end{table}

\textbf{Timer.} Active training time is measured per optimizer step, from the moment the data loader returns the
batch to a CUDA synchronisation after the update (copy to the GPU, forward and backward pass, optimizer step).
Waiting for data (logged separately), logging, validation, checkpoint writes and start-up are excluded; after a crash
the clock resumes from the last checkpoint.

\textbf{Storage.} The shared file system delivers at most $\sim$98 samples per second into one node
(Fig.~\ref{fig:loader}) and stalls for 5--22 minutes at a time on all nodes; TFNO and AViT consume 116 and 330 samples
per second. For the 100\% runs, each job copied the training split (294\,GB) to the GPU node's local disk (1.83\,TB)
and deleted it at exit. Write stalls of the shared file system still hit the checkpoint saves of the 100\% runs, up to
22\% of their wall time, which explains their larger ``other'' time.

\begin{figure}[!htbp]
\centering
\figureorpending{fig_loader.pdf}{scale=1}
\caption{Training samples per second read from the shared file system by one node, against the rate at which each
model consumes them on the GPU (horizontal lines).}
\label{fig:loader}
\end{figure}

\section{Encountered Issues}
\label{app:issues}

\begin{table}[!htbp]
\centering\footnotesize
\begin{tabularx}{\linewidth}{@{}>{\raggedright\arraybackslash}p{0.13\linewidth}>{\raggedright\arraybackslash}p{0.4\linewidth}>{\raggedright\arraybackslash}X@{}}
\toprule
Area & Issue & Effect and handling \\
\midrule
Data on HPC3 & The shared file system delivers at most $\sim$98 training samples per second into one node and stalls
for 5--22 minutes at a time & 10\% runs waited 0.8--7.3\,h for data per 4-hour run (Fig.~\ref{fig:time}); the 100\%
runs copied the data to the node's disk and read errors are retried \\ \addlinespace[3pt]
\midrule
Data on HPC3 & Input/output errors and launch failures of the shared file system & Two 10\% AViT runs crashed and
resumed from their checkpoints; one 100\% run restarted inside its job; six jobs failed at launch and were
resubmitted \\ \addlinespace[3pt]
\midrule
Compute on HPC3 & Queue waits on the paid partition & Median 1--8\,h by purpose; up to 12--15\,h \\
\addlinespace[3pt]
\midrule
Training & TFNO's Fourier layers fail in bf16 & Every model trains in fp32, which favours TFNO \\ \addlinespace[3pt]
\midrule
Data & The files store an evenly spaced placeholder for $y$ & Identified the Chebyshev--Gauss nodes from the data
(Appendix~\ref{app:rb_regimes}) \\ \addlinespace[3pt]
\midrule
Upstream code & The released TFNO was trained without its Fourier layers & Scored with its paper-era forward pass;
not used as a reference (Appendix~\ref{app:tfno}) \\ \addlinespace[3pt]
\midrule
Solver & Our CFL safety factor (0.5) is far larger than the data generator's (0.0039) & The $\mathrm{Ra}=10^{10}$
comparison is reported with this limitation (Section~\ref{sec:solver}) \\
\bottomrule
\end{tabularx}
\caption{Issues met during the project and their effect.}
\label{tab:issues}
\end{table}

\end{document}
